\RequirePackage{fix-cm}
\documentclass[aspectratio=169,11pt]{beamer}

\usepackage[T1]{fontenc}
\usepackage[brazilian]{babel}
\usepackage{helvet}
\usepackage{graphicx}
\usepackage{tikz}
\usepackage{textcomp}

\setbeamertemplate{navigation symbols}{}
\hypersetup{hidelinks,pdfencoding=auto}
\usefonttheme{professionalfonts}
\title{Introdução à Estatística}
\subtitle{Minicurso com Python -- EMASC 2026}
\author{Gabriel C. Souza}
\institute{Liga Acadêmica Matemática \& Empresa -- UFS}
\date{2026}

% Reconstrução a partir do PDF de Gabriel; não é o código-fonte original.
% Cada página de fundos.pdf conserva os desenhos, imagens e equações do PDF.
% Os textos recuperados são posicionados em coordenadas da página (bp).
\newcommand{\Fundo}[1]{%
  \node[anchor=north west,inner sep=0pt] at (0,0)
    {\includegraphics[page=#1,width=\paperwidth,height=\paperheight]{figuras/fundos.pdf}};%
}

\newcommand{\fonter}{\fontfamily{phv}\fontseries{m}\fontshape{n}\selectfont}
\newcommand{\fonteb}{\fontfamily{phv}\fontseries{b}\fontshape{n}\selectfont}
\newcommand{\fontei}{\fontfamily{phv}\fontseries{m}\fontshape{sl}\selectfont}
\newcommand{\fontet}{\fontfamily{cmtt}\fontseries{m}\fontshape{n}\selectfont}
\newcommand{\fontec}{\fontfamily{cmtt}\fontseries{m}\fontshape{it}\selectfont}

\newsavebox{\caixatexto}
% Uso: \Texto{x}{y}{largura}{tamanho}{fonte r/b/i/t/c}{cor HTML}{texto}.
% x,y são contados a partir do canto superior esquerdo, y é a linha de base.
% A largura preserva o espaço originalmente ocupado pelo texto.
\newcommand{\Texto}[7]{%
  \node[anchor=base west,inner sep=0pt,outer sep=0pt] at (#1,#2)
    {\sbox{\caixatexto}{\normalfont\fontsize{#4bp}{#4bp}\selectfont\csname fonte#5\endcsname\color[HTML]{#6}#7}%
     \pgfmathparse{(#3bp)/(\the\wd\caixatexto)}%
     \scalebox{\pgfmathresult}[1]{\usebox{\caixatexto}}};%
}

% Conserva os destinos clicáveis do roteiro original.
\newcommand{\Link}[5]{%
  \node[anchor=north west,inner sep=0pt,outer sep=0pt] at (#1,#2)
    {\hyperlink{#5}{\phantom{\rule{#3bp}{#4bp}}}};%
}

\begin{document}

% Página 1 do PDF de Gabriel
\begin{frame}[plain,noframenumbering,label=original-01]
  \begin{tikzpicture}[remember picture,overlay,x=1bp,y=-1bp]
  \begin{scope}[shift={(current page.north west)}]
  \Fundo{1}
  \Texto{44.599}{25.392}{45.445}{7.572}{b}{0E103E}{EMASC 2026}
  \Texto{108.235}{26.134}{164.902}{7.572}{b}{34E0C8}{LIGA ACADÊMICA MATEMÁTICA \& EMPRESA}
  \Texto{38.787}{55.914}{228.958}{20.822}{b}{FFFFFF}{Introdução à Estatística}
  \Texto{38.787}{140.579}{239.314}{13.25}{b}{34E0C8}{Mini-curso com Python \textbullet{} EMASC 2026}
  \Texto{38.787}{163.983}{97.778}{12.304}{b}{FFFFFF}{Gabriel C. Souza}
  \Texto{38.787}{175.098}{143.069}{8.045}{r}{D1CEE4}{Liga Acadêmica Matemática \& Empresa}
  \Texto{38.787}{186.056}{142.103}{8.045}{r}{D1CEE4}{Universidade Federal de Sergipe \textendash{} UFS}
  \Texto{38.787}{201.822}{90.481}{7.572}{b}{FFB014}{São Cristóvão \textendash{} SE, 2026}
  \end{scope}
  \end{tikzpicture}
\end{frame}


% Página 2 do PDF de Gabriel
\begin{frame}[plain,noframenumbering,label=original-02]
  \begin{tikzpicture}[remember picture,overlay,x=1bp,y=-1bp]
  \begin{scope}[shift={(current page.north west)}]
  \Fundo{2}
  \Texto{40.175}{29.691}{62.579}{17.983}{b}{2C2084}{Roteiro}
  \Texto{40.86}{55.887}{4.736}{8.518}{b}{FFFFFF}{1}
  \Texto{57.649}{55.887}{99.445}{11.357}{r}{2C2084}{Análise exploratória}
  \Texto{40.86}{79.657}{4.736}{8.518}{b}{FFFFFF}{2}
  \Texto{57.649}{79.657}{107.327}{11.357}{r}{2C2084}{Qualidade dos dados}
  \Texto{40.86}{103.427}{4.736}{8.518}{b}{FFFFFF}{3}
  \Texto{57.649}{103.427}{106.124}{11.357}{r}{2C2084}{Perguntas da análise}
  \Texto{40.86}{127.25}{4.736}{8.518}{b}{FFFFFF}{4}
  \Texto{57.649}{127.25}{88.247}{11.357}{r}{2C2084}{Análise descritiva}
  \Texto{40.86}{151.02}{4.736}{8.518}{b}{FFFFFF}{5}
  \Texto{57.649}{151.02}{64.942}{11.357}{r}{2C2084}{Distribuições}
  \Texto{40.86}{174.843}{4.736}{8.518}{b}{FFFFFF}{6}
  \Texto{57.649}{174.843}{86.464}{11.357}{r}{2C2084}{Correlação linear}
  \Texto{40.86}{198.666}{4.736}{8.518}{b}{FFFFFF}{7}
  \Texto{57.649}{198.666}{69.439}{11.357}{r}{2C2084}{Probabilidade}
  \Texto{40.86}{222.436}{4.736}{8.518}{b}{FFFFFF}{8}
  \Texto{57.649}{222.436}{93.437}{11.357}{r}{2C2084}{Síntese da análise}
  \Texto{37.209}{247.463}{40.499}{6.625}{b}{737696}{EMASC 2026}
  \Texto{80.653}{247.463}{2.319}{6.625}{r}{FFB014}{\textbullet{}}
  \Texto{85.913}{247.463}{117.821}{6.625}{r}{737696}{Liga Acadêmica Matemática \& Empresa}
  \Texto{206.682}{247.463}{2.319}{6.625}{r}{FFB014}{\textbullet{}}
  \Texto{211.943}{247.463}{13.25}{6.625}{r}{737696}{UFS}
  \Texto{401.232}{248.132}{3.684}{6.625}{b}{2C2084}{1}
  \Texto{406.02}{248.132}{10.315}{6.625}{r}{737696}{/ 40}
  \Link{56.653}{44.08}{101.436}{17.584}{original-03}
  \Link{56.653}{67.85}{109.318}{17.583}{original-08}
  \Link{56.653}{91.673}{108.114}{17.583}{original-14}
  \Link{56.653}{115.443}{90.238}{17.583}{original-17}
  \Link{56.653}{139.266}{66.933}{17.583}{original-26}
  \Link{56.653}{163.089}{88.455}{17.583}{original-37}
  \Link{56.653}{186.859}{71.43}{17.583}{original-52}
  \Link{56.653}{210.682}{95.428}{17.583}{original-63}
  \end{scope}
  \end{tikzpicture}
\end{frame}


\section{Análise exploratória}

% Página 3 do PDF de Gabriel
\begin{frame}[plain,noframenumbering,label=original-03]
  \begin{tikzpicture}[remember picture,overlay,x=1bp,y=-1bp]
  \begin{scope}[shift={(current page.north west)}]
  \Fundo{3}
  \Texto{44.456}{41.894}{106.322}{8.518}{b}{34E0C8}{ANÁLISE EXPLORATÓRIA}
  \Texto{43.039}{116.125}{63.147}{56.787}{b}{FFB014}{01}
  \Texto{43.039}{155.148}{235.888}{24.608}{b}{FFFFFF}{Análise exploratória}
  \Texto{43.039}{234.662}{46.285}{7.572}{b}{34E0C8}{EMASC 2026}
  \Texto{91.429}{234.662}{2.65}{7.572}{b}{FFB014}{\textbullet{}}
  \Texto{96.176}{234.662}{21.45}{7.572}{b}{34E0C8}{LAME}
  \Link{42.042}{135.586}{237.879}{25.652}{original-03}
  \end{scope}
  \end{tikzpicture}
\end{frame}


% Página 4 do PDF de Gabriel
\begin{frame}[plain,noframenumbering,label=original-04]
  \begin{tikzpicture}[remember picture,overlay,x=1bp,y=-1bp]
  \begin{scope}[shift={(current page.north west)}]
  \Fundo{4}
  \Texto{40.175}{28.167}{238.844}{17.983}{b}{2C2084}{Etapas da análise de dados}
  \Texto{35.433}{63.506}{274.953}{8.518}{i}{737696}{Analisar dados é um processo em etapas: cada uma prepara a seguinte.}
  \Texto{28.574}{100.879}{55.225}{9.465}{b}{027499}{Exploratória}
  \Texto{28.574}{115.823}{48.922}{9.465}{r}{000000}{Conhecer a}
  \Texto{28.574}{127.778}{40.025}{9.465}{r}{000000}{estrutura,}
  \Texto{28.574}{139.733}{40.925}{9.465}{r}{000000}{verificar a}
  \Texto{28.574}{151.689}{48.931}{9.465}{r}{000000}{qualidade e}
  \Texto{28.574}{163.644}{40.499}{9.465}{r}{000000}{identificar}
  \Texto{28.574}{175.599}{36.684}{9.465}{r}{000000}{padrões.}
  \Texto{137.268}{100.643}{45.051}{9.465}{b}{007078}{Descritiva}
  \Texto{137.268}{115.586}{47.853}{9.465}{r}{000000}{Resumir os}
  \Texto{137.268}{127.542}{41.559}{9.465}{r}{000000}{dados em}
  \Texto{137.268}{139.497}{43.66}{9.465}{r}{000000}{medidas e}
  \Texto{137.268}{151.452}{35.53}{9.465}{r}{000000}{gráficos.}
  \Texto{245.963}{100.643}{39.789}{9.465}{b}{231A6A}{Preditiva}
  \Texto{245.963}{115.586}{64.974}{9.465}{r}{000000}{Estimar valores}
  \Texto{245.963}{127.542}{67.387}{9.465}{r}{000000}{não observados}
  \Texto{245.963}{139.497}{49.973}{9.465}{r}{000000}{por meio de}
  \Texto{245.963}{151.452}{38.255}{9.465}{r}{000000}{modelos.}
  \Texto{354.657}{100.643}{48.212}{9.465}{b}{CC8D10}{Prescritiva}
  \Texto{354.657}{115.586}{54.175}{9.465}{r}{000000}{Recomendar}
  \Texto{354.657}{127.542}{57.715}{9.465}{r}{000000}{ações a partir}
  \Texto{354.657}{139.497}{67.472}{9.465}{r}{000000}{das estimativas.}
  \Texto{49.227}{210.138}{357.083}{8.518}{r}{000000}{As duas primeiras etapas descrevem o que os dados mostram. As duas últimas dependem de}
  \Texto{49.227}{221.097}{157.796}{8.518}{r}{000000}{modelos ajustados e de dados confiáveis.}
  \Texto{37.21}{247.463}{40.499}{6.625}{b}{737696}{EMASC 2026}
  \Texto{80.653}{247.463}{2.319}{6.625}{r}{FFB014}{\textbullet{}}
  \Texto{85.914}{247.463}{117.821}{6.625}{r}{737696}{Liga Acadêmica Matemática \& Empresa}
  \Texto{206.683}{247.463}{2.319}{6.625}{r}{FFB014}{\textbullet{}}
  \Texto{211.943}{247.463}{13.25}{6.625}{r}{737696}{UFS}
  \Texto{401.232}{248.132}{3.684}{6.625}{b}{2C2084}{2}
  \Texto{406.02}{248.132}{10.315}{6.625}{r}{737696}{/ 40}
  \end{scope}
  \end{tikzpicture}
\end{frame}


% Página 5 do PDF de Gabriel
\begin{frame}[plain,noframenumbering,label=original-05]
  \begin{tikzpicture}[remember picture,overlay,x=1bp,y=-1bp]
  \begin{scope}[shift={(current page.north west)}]
  \Fundo{5}
  \Texto{40.175}{27.942}{181.839}{17.983}{b}{2C2084}{O conjunto de dados}
  \Texto{35.433}{63.506}{214.773}{8.518}{i}{737696}{Toda análise começa por conhecer o material disponível.}
  \Texto{34.223}{100.643}{99.008}{9.465}{b}{027499}{Unidade de observa-}
  \Texto{34.223}{112.598}{16.307}{9.465}{b}{027499}{ção}
  \Texto{34.223}{127.542}{56.806}{9.465}{r}{000000}{Cada registro}
  \Texto{34.223}{139.497}{60.497}{9.465}{r}{000000}{corresponde a}
  \Texto{97.35}{139.497}{14.197}{9.465}{b}{000000}{um}
  \Texto{34.223}{151.452}{44.701}{9.465}{b}{000000}{estudante}
  \Texto{78.923}{151.452}{2.631}{9.465}{r}{000000}{.}
  \Texto{177.266}{100.643}{43.461}{9.465}{b}{007078}{Registros}
  \Texto{177.266}{115.586}{74.457}{9.465}{r}{000000}{O arquivo contém}
  \Texto{177.266}{127.542}{58.926}{9.465}{b}{000000}{16.000 linhas}
  \Texto{238.822}{127.542}{18.418}{9.465}{r}{000000}{e 10}
  \Texto{177.266}{139.497}{35.104}{9.465}{r}{000000}{colunas.}
  \Texto{320.309}{100.879}{44.161}{9.465}{b}{CC8D10}{Conteúdo}
  \Texto{320.309}{115.823}{65.22}{9.465}{r}{000000}{Hábitos diários,}
  \Texto{320.309}{127.778}{86.789}{9.465}{r}{000000}{bem-estar e perfil do}
  \Texto{320.309}{139.733}{44.057}{9.465}{r}{000000}{estudante.}
  \Texto{49.227}{184.118}{145.045}{8.518}{r}{000000}{As variáveis de hábito são medidas em}
  \Texto{196.169}{184.118}{53.016}{8.518}{b}{000000}{horas por dia}
  \Texto{249.185}{184.118}{157.123}{8.518}{r}{000000}{. As pontuações de bem-estar são índices}
  \Texto{49.227}{195.076}{254.697}{8.518}{r}{000000}{sem unidade de medida: valores maiores indicam melhor condição.}
  \Texto{37.21}{247.463}{40.499}{6.625}{b}{737696}{EMASC 2026}
  \Texto{80.653}{247.463}{2.319}{6.625}{r}{FFB014}{\textbullet{}}
  \Texto{85.913}{247.463}{117.821}{6.625}{r}{737696}{Liga Acadêmica Matemática \& Empresa}
  \Texto{206.682}{247.463}{2.319}{6.625}{r}{FFB014}{\textbullet{}}
  \Texto{211.943}{247.463}{13.25}{6.625}{r}{737696}{UFS}
  \Texto{401.232}{248.132}{3.684}{6.625}{b}{2C2084}{3}
  \Texto{406.02}{248.132}{10.315}{6.625}{r}{737696}{/ 40}
  \end{scope}
  \end{tikzpicture}
\end{frame}


% Página 6 do PDF de Gabriel
\begin{frame}[plain,noframenumbering,label=original-06]
  \begin{tikzpicture}[remember picture,overlay,x=1bp,y=-1bp]
  \begin{scope}[shift={(current page.north west)}]
  \Fundo{6}
  \Texto{40.175}{29.916}{197.52}{17.983}{b}{2C2084}{Dicionário de variáveis}
  \Texto{76.421}{62.847}{28.59}{7.572}{b}{262848}{Variável}
  \Texto{194.182}{62.847}{36.192}{7.572}{b}{262848}{Descrição}
  \Texto{346.046}{62.847}{15.984}{7.572}{b}{262848}{Tipo}
  \Texto{76.421}{77.457}{8.464}{7.97}{t}{2C2084}{id}
  \Texto{194.18}{77.457}{78.692}{7.572}{r}{262848}{identificador do registro}
  \Texto{346.051}{77.457}{26.508}{7.572}{r}{262848}{nominal}
  \Texto{76.421}{86.922}{21.161}{7.97}{t}{2C2084}{idade}
  \Texto{194.18}{86.922}{49.662}{7.572}{r}{262848}{idade em anos}
  \Texto{346.051}{86.922}{26.508}{7.572}{r}{262848}{discreta}
  \Texto{76.421}{96.386}{25.393}{7.97}{t}{2C2084}{genero}
  \Texto{194.18}{96.386}{58.846}{7.572}{r}{262848}{gênero declarado}
  \Texto{346.051}{96.386}{26.508}{7.572}{r}{262848}{nominal}
  \Texto{76.421}{105.851}{21.161}{7.97}{t}{2C2084}{nivel}
  \Texto{194.18}{105.851}{50.737}{7.572}{r}{262848}{nível de ensino}
  \Texto{346.051}{105.851}{22.722}{7.572}{r}{262848}{ordinal}
  \Texto{76.421}{115.315}{21.161}{7.97}{t}{2C2084}{redes}
  \Texto{97.587}{115.315}{2.105}{7.572}{r}{262848}{,}
  \Texto{101.796}{115.315}{8.464}{7.97}{t}{2C2084}{ia}
  \Texto{194.18}{115.315}{118.17}{7.572}{r}{262848}{horas/dia em redes sociais e em IA}
  \Texto{346.051}{115.315}{28.969}{7.572}{r}{262848}{contínua}
  \Texto{76.421}{124.78}{16.928}{7.97}{t}{2C2084}{sono}
  \Texto{194.18}{124.78}{60.111}{7.572}{r}{262848}{horas/dia de sono}
  \Texto{346.051}{124.78}{28.969}{7.572}{r}{262848}{contínua}
  \Texto{76.421}{134.244}{38.089}{7.97}{t}{2C2084}{atividade}
  \Texto{194.18}{134.244}{93.774}{7.572}{r}{262848}{horas/dia de atividade física}
  \Texto{346.051}{134.244}{28.969}{7.572}{r}{262848}{contínua}
  \Texto{76.421}{143.709}{50.785}{7.97}{t}{2C2084}{saude\_mental}
  \Texto{127.218}{143.709}{2.105}{7.572}{r}{262848}{,}
  \Texto{131.428}{143.709}{50.785}{7.97}{t}{2C2084}{saude\_fisica}
  \Texto{194.18}{143.709}{139.916}{7.572}{r}{262848}{pontuação de bem-estar (maior = melhor)}
  \Texto{346.051}{143.709}{28.969}{7.572}{r}{262848}{contínua}
  \Texto{49.227}{172.851}{357.083}{8.518}{r}{000000}{Cada linha é um estudante e cada coluna, uma variável. O dicionário registra o que cada coluna}
  \Texto{49.227}{183.81}{285.966}{8.518}{r}{000000}{mede e em que unidade, o que determina quais cálculos são interpretáveis.}
  \Texto{37.21}{247.463}{40.499}{6.625}{b}{737696}{EMASC 2026}
  \Texto{80.653}{247.463}{2.319}{6.625}{r}{FFB014}{\textbullet{}}
  \Texto{85.913}{247.463}{117.821}{6.625}{r}{737696}{Liga Acadêmica Matemática \& Empresa}
  \Texto{206.682}{247.463}{2.319}{6.625}{r}{FFB014}{\textbullet{}}
  \Texto{211.943}{247.463}{13.25}{6.625}{r}{737696}{UFS}
  \Texto{401.232}{248.132}{3.684}{6.625}{b}{2C2084}{4}
  \Texto{406.02}{248.132}{10.315}{6.625}{r}{737696}{/ 40}
  \end{scope}
  \end{tikzpicture}
\end{frame}


% Página 7 do PDF de Gabriel
\begin{frame}[plain,noframenumbering,label=original-07]
  \begin{tikzpicture}[remember picture,overlay,x=1bp,y=-1bp]
  \begin{scope}[shift={(current page.north west)}]
  \Fundo{7}
  \Texto{40.175}{28.028}{234.528}{17.983}{b}{2C2084}{Classificação das variáveis}
  \Texto{35.433}{63.506}{213.367}{8.518}{i}{737696}{O tipo da variável define quais operações fazem sentido.}
  \Texto{28.574}{100.643}{37.337}{9.465}{b}{CC8D10}{Nominal}
  \Texto{28.574}{132.523}{31.365}{9.465}{b}{000000}{gênero}
  \Texto{28.574}{144.478}{65.635}{6.625}{r}{000000}{categorias sem ordem}
  \Texto{28.574}{156.433}{21.724}{6.625}{r}{000000}{entre si}
  \Texto{137.268}{100.642}{32.946}{9.465}{b}{027499}{Ordinal}
  \Texto{137.268}{134.438}{21.57}{9.465}{b}{000000}{nível}
  \Texto{137.268}{146.393}{59.745}{6.625}{r}{000000}{categorias com uma}
  \Texto{137.268}{158.348}{40.804}{6.625}{r}{000000}{ordem natural}
  \Texto{245.963}{100.642}{37.347}{9.465}{b}{007078}{Discreta}
  \Texto{245.963}{132.523}{24.721}{9.465}{b}{000000}{idade}
  \Texto{245.963}{144.478}{63.17}{6.625}{r}{000000}{valores que podemos}
  \Texto{245.963}{156.433}{52.557}{6.625}{r}{000000}{contar: 1, 2, 3, . . .}
  \Texto{354.657}{100.879}{40.915}{9.465}{b}{231A6A}{Contínua}
  \Texto{354.657}{132.759}{64.69}{9.465}{b}{000000}{horas de sono}
  \Texto{354.657}{144.715}{49.907}{6.625}{r}{000000}{valores medidos:}
  \Texto{354.657}{156.67}{61.123}{6.625}{r}{000000}{qualquer número em}
  \Texto{354.657}{168.625}{36.113}{6.625}{r}{000000}{um intervalo}
  \Texto{49.227}{198.414}{357.083}{8.518}{r}{000000}{Classificar evita operações sem sentido: somar categorias não produz informação; somar ho-}
  \Texto{49.227}{209.373}{31.968}{8.518}{r}{000000}{ras, sim.}
  \Texto{37.21}{247.463}{40.499}{6.625}{b}{737696}{EMASC 2026}
  \Texto{80.653}{247.463}{2.319}{6.625}{r}{FFB014}{\textbullet{}}
  \Texto{85.914}{247.463}{117.821}{6.625}{r}{737696}{Liga Acadêmica Matemática \& Empresa}
  \Texto{206.683}{247.463}{2.319}{6.625}{r}{FFB014}{\textbullet{}}
  \Texto{211.943}{247.463}{13.25}{6.625}{r}{737696}{UFS}
  \Texto{401.232}{248.132}{3.684}{6.625}{b}{2C2084}{5}
  \Texto{406.02}{248.132}{10.315}{6.625}{r}{737696}{/ 40}
  \end{scope}
  \end{tikzpicture}
\end{frame}


\section{Qualidade dos dados}

% Página 8 do PDF de Gabriel
\begin{frame}[plain,noframenumbering,label=original-08]
  \begin{tikzpicture}[remember picture,overlay,x=1bp,y=-1bp]
  \begin{scope}[shift={(current page.north west)}]
  \Fundo{8}
  \Texto{44.456}{41.894}{106.322}{8.518}{b}{34E0C8}{ANÁLISE EXPLORATÓRIA}
  \Texto{43.039}{116.125}{63.147}{56.787}{b}{FFB014}{02}
  \Texto{43.039}{154.533}{250.948}{24.608}{b}{FFFFFF}{Qualidade dos dados}
  \Texto{43.039}{234.662}{46.285}{7.572}{b}{34E0C8}{EMASC 2026}
  \Texto{91.429}{234.662}{2.65}{7.572}{b}{FFB014}{\textbullet{}}
  \Texto{96.176}{234.662}{21.45}{7.572}{b}{34E0C8}{LAME}
  \Link{42.042}{135.586}{252.94}{21.223}{original-08}
  \end{scope}
  \end{tikzpicture}
\end{frame}


% Página 9 do PDF de Gabriel
\begin{frame}[plain,noframenumbering,label=original-09]
  \begin{tikzpicture}[remember picture,overlay,x=1bp,y=-1bp]
  \begin{scope}[shift={(current page.north west)}]
  \Fundo{9}
  \Texto{40.175}{27.803}{210.485}{17.983}{b}{2C2084}{Protocolo de verificação}
  \Texto{35.433}{63.506}{297.764}{8.518}{i}{737696}{Erros nos dados produzem conclusões erradas, mesmo com cálculos corretos.}
  \Texto{28.574}{100.439}{26.292}{9.465}{b}{027499}{Nulos}
  \Texto{28.574}{115.383}{44.18}{9.465}{r}{000000}{Há células}
  \Texto{28.574}{127.338}{31.848}{9.465}{r}{000000}{vazias?}
  \Texto{137.268}{100.642}{47.862}{9.465}{b}{CC8D10}{Duplicatas}
  \Texto{137.268}{115.586}{39.448}{9.465}{r}{000000}{Há linhas}
  \Texto{137.268}{127.542}{44.19}{9.465}{r}{000000}{repetidas?}
  \Texto{245.963}{100.642}{29.274}{9.465}{b}{007078}{Faixas}
  \Texto{245.963}{115.586}{62.882}{9.465}{r}{000000}{Os valores são}
  \Texto{245.963}{127.542}{47.105}{9.465}{r}{000000}{plausíveis?}
  \Texto{354.657}{100.879}{58.907}{9.465}{b}{231A6A}{Consistência}
  \Texto{354.657}{115.823}{50.493}{9.465}{r}{000000}{Uma coluna}
  \Texto{354.657}{127.778}{49.992}{9.465}{r}{000000}{depende de}
  \Texto{354.657}{139.733}{26.737}{9.465}{r}{000000}{outra?}
  \Texto{49.227}{169.555}{357.083}{8.518}{r}{000000}{Nenhuma estatística é mais confiável que os dados que a originam. As quatro verificações}
  \Texto{49.227}{180.514}{108.553}{8.518}{r}{000000}{antecedem qualquer cálculo.}
  \Texto{37.21}{247.463}{40.499}{6.625}{b}{737696}{EMASC 2026}
  \Texto{80.653}{247.463}{2.319}{6.625}{r}{FFB014}{\textbullet{}}
  \Texto{85.914}{247.463}{117.821}{6.625}{r}{737696}{Liga Acadêmica Matemática \& Empresa}
  \Texto{206.683}{247.463}{2.319}{6.625}{r}{FFB014}{\textbullet{}}
  \Texto{211.943}{247.463}{13.25}{6.625}{r}{737696}{UFS}
  \Texto{401.232}{248.132}{3.684}{6.625}{b}{2C2084}{6}
  \Texto{406.02}{248.132}{10.315}{6.625}{r}{737696}{/ 40}
  \end{scope}
  \end{tikzpicture}
\end{frame}


% Página 10 do PDF de Gabriel
\begin{frame}[plain,noframenumbering,label=original-10]
  \begin{tikzpicture}[remember picture,overlay,x=1bp,y=-1bp]
  \begin{scope}[shift={(current page.north west)}]
  \Fundo{10}
  \Texto{40.175}{27.803}{183.493}{17.983}{b}{2C2084}{Registros duplicados}
  \Texto{35.433}{63.506}{268.13}{8.518}{i}{737696}{Uma mesma observação contada duas vezes distorce qualquer média.}
  \Texto{49.227}{203.827}{60.861}{8.518}{b}{000000}{1.000 registros}
  \Texto{113.365}{203.827}{292.943}{8.518}{r}{000000}{repetem integralmente observações anteriores e diferem apenas no identifi-}
  \Texto{49.227}{214.786}{327.5}{8.518}{r}{000000}{cador. Mantidos, receberiam peso duplo nas estimativas. Após a remoção, restam}
  \Texto{380.261}{214.786}{26.048}{8.518}{b}{000000}{15.000}
  \Texto{49.227}{225.745}{66.1}{8.518}{b}{000000}{registros únicos}
  \Texto{115.326}{225.745}{2.368}{8.518}{r}{000000}{.}
  \Texto{37.21}{247.463}{40.499}{6.625}{b}{737696}{EMASC 2026}
  \Texto{80.653}{247.463}{2.319}{6.625}{r}{FFB014}{\textbullet{}}
  \Texto{85.913}{247.463}{117.821}{6.625}{r}{737696}{Liga Acadêmica Matemática \& Empresa}
  \Texto{206.682}{247.463}{2.319}{6.625}{r}{FFB014}{\textbullet{}}
  \Texto{211.943}{247.463}{13.25}{6.625}{r}{737696}{UFS}
  \Texto{401.232}{248.132}{3.684}{6.625}{b}{2C2084}{7}
  \Texto{406.02}{248.132}{10.315}{6.625}{r}{737696}{/ 40}
  \end{scope}
  \end{tikzpicture}
\end{frame}


% Página 11 do PDF de Gabriel
\begin{frame}[plain,noframenumbering,label=original-11]
  \begin{tikzpicture}[remember picture,overlay,x=1bp,y=-1bp]
  \begin{scope}[shift={(current page.north west)}]
  \Fundo{11}
  \Texto{40.175}{29.916}{156.34}{17.983}{b}{2C2084}{Acúmulo de zeros}
  \Texto{35.433}{63.506}{319.519}{8.518}{i}{737696}{Muitos valores repetidos em um ponto exato podem indicar um problema de registro.}
  \Texto{49.227}{203.827}{324.706}{8.518}{r}{000000}{As proporções de zeros exatos (redes 3,6\%; IA 8,3\%; atividade 17,6\%) sugerem que}
  \Texto{376.65}{203.827}{29.66}{8.518}{b}{000000}{valores}
  \Texto{49.227}{214.786}{150.317}{8.518}{b}{000000}{negativos foram convertidos em zero}
  \Texto{199.542}{214.786}{205.003}{8.518}{r}{000000}{. Esse ajuste desloca a média e a mediana para cima.}
  \Texto{37.21}{247.463}{40.499}{6.625}{b}{737696}{EMASC 2026}
  \Texto{80.653}{247.463}{2.319}{6.625}{r}{FFB014}{\textbullet{}}
  \Texto{85.913}{247.463}{117.821}{6.625}{r}{737696}{Liga Acadêmica Matemática \& Empresa}
  \Texto{206.682}{247.463}{2.319}{6.625}{r}{FFB014}{\textbullet{}}
  \Texto{211.943}{247.463}{13.25}{6.625}{r}{737696}{UFS}
  \Texto{401.232}{248.132}{3.684}{6.625}{b}{2C2084}{8}
  \Texto{406.02}{248.132}{10.315}{6.625}{r}{737696}{/ 40}
  \end{scope}
  \end{tikzpicture}
\end{frame}


% Página 12 do PDF de Gabriel
\begin{frame}[plain,noframenumbering,label=original-12]
  \begin{tikzpicture}[remember picture,overlay,x=1bp,y=-1bp]
  \begin{scope}[shift={(current page.north west)}]
  \Fundo{12}
  \Texto{40.175}{28.167}{244.49}{17.983}{b}{2C2084}{Dependência entre variáveis}
  \Texto{35.433}{63.506}{289.425}{8.518}{i}{737696}{Duas colunas podem parecer diferentes, mas carregar a mesma informação.}
  \Texto{125.333}{87.736}{22.249}{8.518}{b}{2C2084}{Idade}
  \Texto{280.961}{87.934}{62.488}{8.518}{b}{2C2084}{Nível de ensino}
  \Texto{111.125}{113.248}{50.665}{8.518}{b}{262848}{13 a 17 anos}
  \Texto{111.125}{144.429}{50.665}{8.518}{b}{262848}{18 a 21 anos}
  \Texto{111.125}{175.61}{50.665}{8.518}{b}{262848}{22 a 25 anos}
  \Texto{287.401}{112.445}{49.609}{8.518}{b}{262848}{High School}
  \Texto{296.779}{143.626}{30.852}{8.518}{b}{262848}{College}
  \Texto{291.677}{174.873}{41.057}{8.518}{b}{262848}{University}
  \Texto{49.227}{202.872}{334.212}{8.518}{r}{000000}{Cada faixa de idade corresponde a exatamente um nível. Comparar níveis é, na prática,}
  \Texto{385.96}{202.872}{20.35}{8.518}{b}{000000}{com-}
  \Texto{49.227}{213.831}{77.65}{8.518}{b}{000000}{parar faixas etárias}
  \Texto{126.876}{213.831}{193.018}{8.518}{r}{000000}{: as duas variáveis carregam a mesma informação.}
  \Texto{37.21}{247.463}{40.499}{6.625}{b}{737696}{EMASC 2026}
  \Texto{80.653}{247.463}{2.319}{6.625}{r}{FFB014}{\textbullet{}}
  \Texto{85.913}{247.463}{117.821}{6.625}{r}{737696}{Liga Acadêmica Matemática \& Empresa}
  \Texto{206.682}{247.463}{2.319}{6.625}{r}{FFB014}{\textbullet{}}
  \Texto{211.943}{247.463}{13.25}{6.625}{r}{737696}{UFS}
  \Texto{401.232}{248.132}{3.684}{6.625}{b}{2C2084}{9}
  \Texto{406.02}{248.132}{10.315}{6.625}{r}{737696}{/ 40}
  \end{scope}
  \end{tikzpicture}
\end{frame}


% Página 13 do PDF de Gabriel
\begin{frame}[plain,noframenumbering,label=original-13]
  \begin{tikzpicture}[remember picture,overlay,x=1bp,y=-1bp]
  \begin{scope}[shift={(current page.north west)}]
  \Fundo{13}
  \Texto{40.175}{28.028}{178.53}{17.983}{b}{2C2084}{Síntese da validação}
  \Texto{35.433}{63.506}{206.877}{8.518}{i}{737696}{As verificações determinam o alcance das conclusões.}
  \Texto{87.929}{81.728}{45.026}{8.518}{b}{262848}{Verificação}
  \Texto{148.867}{81.728}{41.167}{8.518}{b}{262848}{Resultado}
  \Texto{252.463}{81.728}{58.221}{8.518}{b}{262848}{Consequência}
  \Texto{87.929}{97.833}{21.772}{8.518}{r}{262848}{Nulos}
  \Texto{148.867}{97.833}{30.767}{8.518}{r}{262848}{nenhum}
  \Texto{252.463}{97.833}{8.518}{8.518}{r}{262848}{\textemdash{}}
  \Texto{87.929}{108.792}{39.762}{8.518}{r}{262848}{Duplicatas}
  \Texto{148.867}{108.792}{62.829}{8.518}{r}{262848}{1.000 removidas}
  \Texto{252.463}{108.792}{53.033}{8.518}{r}{262848}{15.000 únicos}
  \Texto{87.929}{119.751}{48.979}{8.518}{r}{262848}{Zeros exatos}
  \Texto{148.867}{119.751}{38.348}{8.518}{r}{262848}{até 17,6\%}
  \Texto{252.463}{119.751}{110.785}{8.518}{r}{262848}{média e mediana deslocadas}
  \Texto{87.929}{130.71}{48.757}{8.518}{r}{262848}{Consistência}
  \Texto{148.867}{130.71}{91.637}{8.518}{r}{262848}{nível definido pela idade}
  \Texto{252.463}{130.71}{93.093}{8.518}{r}{262848}{grupos não comparáveis}
  \Texto{49.227}{160.3}{357.083}{8.518}{r}{000000}{A idade quase uniforme, a ausência de associação com a idade e as cópias ao final do arquivo}
  \Texto{49.227}{171.259}{29.345}{8.518}{r}{000000}{indicam}
  \Texto{80.939}{171.259}{69.413}{8.518}{b}{000000}{dados simulados}
  \Texto{150.352}{171.259}{254.944}{8.518}{r}{000000}{. As conclusões descrevem este conjunto, não uma população real.}
  \Texto{37.21}{247.463}{40.499}{6.625}{b}{737696}{EMASC 2026}
  \Texto{80.653}{247.463}{2.319}{6.625}{r}{FFB014}{\textbullet{}}
  \Texto{85.913}{247.463}{117.821}{6.625}{r}{737696}{Liga Acadêmica Matemática \& Empresa}
  \Texto{206.682}{247.463}{2.319}{6.625}{r}{FFB014}{\textbullet{}}
  \Texto{211.943}{247.463}{13.25}{6.625}{r}{737696}{UFS}
  \Texto{397.549}{248.132}{7.367}{6.625}{b}{2C2084}{10}
  \Texto{406.02}{248.132}{10.315}{6.625}{r}{737696}{/ 40}
  \end{scope}
  \end{tikzpicture}
\end{frame}


\section{Perguntas da análise}

% Página 14 do PDF de Gabriel
\begin{frame}[plain,noframenumbering,label=original-14]
  \begin{tikzpicture}[remember picture,overlay,x=1bp,y=-1bp]
  \begin{scope}[shift={(current page.north west)}]
  \Fundo{14}
  \Texto{44.456}{41.894}{106.322}{8.518}{b}{34E0C8}{ANÁLISE EXPLORATÓRIA}
  \Texto{43.039}{116.125}{63.147}{56.787}{b}{FFB014}{03}
  \Texto{43.039}{155.148}{247.134}{24.608}{b}{FFFFFF}{Perguntas da análise}
  \Texto{43.039}{234.662}{46.285}{7.572}{b}{34E0C8}{EMASC 2026}
  \Texto{91.429}{234.662}{2.65}{7.572}{b}{FFB014}{\textbullet{}}
  \Texto{96.176}{234.662}{21.45}{7.572}{b}{34E0C8}{LAME}
  \Link{42.042}{135.586}{249.126}{26.033}{original-14}
  \end{scope}
  \end{tikzpicture}
\end{frame}


% Página 15 do PDF de Gabriel
\begin{frame}[plain,noframenumbering,label=original-15]
  \begin{tikzpicture}[remember picture,overlay,x=1bp,y=-1bp]
  \begin{scope}[shift={(current page.north west)}]
  \Fundo{15}
  \Texto{40.175}{28.028}{181.084}{17.983}{b}{2C2084}{Perguntas da análise}
  \Texto{35.433}{63.506}{323.139}{8.518}{i}{737696}{A exploração revelou padrões; agora é preciso transformá-los em perguntas precisas.}
  \Texto{34.223}{100.879}{7.893}{9.465}{b}{027499}{1.}
  \Texto{50.142}{100.879}{83.089}{9.465}{b}{027499}{Valor típico e dis-}
  \Texto{34.223}{112.834}{30.892}{9.465}{b}{027499}{persão}
  \Texto{34.223}{127.778}{83.931}{9.465}{r}{000000}{Qual é o valor típico}
  \Texto{34.223}{139.733}{74.183}{9.465}{r}{000000}{dos hábitos e das}
  \Texto{34.223}{151.689}{82.994}{9.465}{r}{000000}{pontuações, e quão}
  \Texto{34.223}{163.644}{63.649}{9.465}{r}{000000}{dispersos são?}
  \Texto{177.266}{100.643}{63.753}{9.465}{b}{007078}{2. Associação}
  \Texto{177.266}{115.586}{96.263}{9.465}{r}{000000}{O uso de redes sociais}
  \Texto{177.266}{127.542}{76.284}{9.465}{r}{000000}{está associado ao}
  \Texto{177.266}{139.497}{83.732}{9.465}{r}{000000}{bem-estar? Em que}
  \Texto{177.266}{151.452}{95.222}{9.465}{r}{000000}{direção e intensidade?}
  \Texto{320.309}{100.643}{97.267}{9.465}{b}{CC8D10}{3. Chance de eventos}
  \Texto{320.309}{115.586}{87.329}{9.465}{r}{000000}{Qual a probabilidade}
  \Texto{320.309}{127.542}{70.88}{9.465}{r}{000000}{de sono curto ou}
  \Texto{320.309}{139.497}{74.713}{9.465}{r}{000000}{pontuação baixa?}
  \Texto{320.309}{151.452}{62.059}{9.465}{r}{000000}{Como modelar}
  \Texto{320.309}{163.407}{48.931}{9.465}{r}{000000}{contagens?}
  \Texto{49.227}{198.345}{178.427}{8.518}{r}{000000}{Uma pergunta respondível por dados identifica}
  \Texto{230.43}{198.345}{124.056}{8.518}{b}{000000}{a variável, o grupo e a medida}
  \Texto{357.263}{198.345}{49.047}{8.518}{r}{000000}{de interesse.}
  \Texto{49.227}{209.304}{173.98}{8.518}{r}{000000}{Perguntas vagas não se traduzem em cálculo.}
  \Texto{37.21}{247.463}{40.499}{6.625}{b}{737696}{EMASC 2026}
  \Texto{80.653}{247.463}{2.319}{6.625}{r}{FFB014}{\textbullet{}}
  \Texto{85.913}{247.463}{117.821}{6.625}{r}{737696}{Liga Acadêmica Matemática \& Empresa}
  \Texto{206.682}{247.463}{2.319}{6.625}{r}{FFB014}{\textbullet{}}
  \Texto{211.943}{247.463}{13.25}{6.625}{r}{737696}{UFS}
  \Texto{397.549}{248.132}{7.367}{6.625}{b}{2C2084}{11}
  \Texto{406.02}{248.132}{10.315}{6.625}{r}{737696}{/ 40}
  \end{scope}
  \end{tikzpicture}
\end{frame}


% Página 16 do PDF de Gabriel
\begin{frame}[plain,noframenumbering,label=original-16]
  \begin{tikzpicture}[remember picture,overlay,x=1bp,y=-1bp]
  \begin{scope}[shift={(current page.north west)}]
  \Fundo{16}
  \Texto{40.175}{28.028}{184.051}{17.983}{b}{2C2084}{Perguntas e métodos}
  \Texto{35.433}{63.506}{210.437}{8.518}{i}{737696}{Cada pergunta exige um método capaz de respondê-la.}
  \Texto{99.152}{81.728}{36.534}{8.518}{b}{262848}{Pergunta}
  \Texto{192.91}{81.728}{30.281}{8.518}{b}{262848}{Método}
  \Texto{326.932}{81.728}{25.094}{8.518}{b}{262848}{Seção}
  \Texto{99.152}{97.833}{42.011}{8.518}{r}{262848}{Valor típico}
  \Texto{192.91}{97.833}{86.637}{8.518}{r}{262848}{média, mediana, moda}
  \Texto{337.111}{97.833}{4.736}{8.518}{r}{262848}{4}
  \Texto{99.152}{108.792}{38.34}{8.518}{r}{262848}{Dispersão}
  \Texto{192.91}{108.792}{111.773}{8.518}{r}{262848}{variância, desvio padrão, IQR}
  \Texto{337.111}{108.792}{4.736}{8.518}{r}{262848}{4}
  \Texto{99.152}{119.751}{81.807}{8.518}{r}{262848}{Forma da distribuição}
  \Texto{192.91}{119.751}{73.902}{8.518}{r}{262848}{histograma, boxplot}
  \Texto{337.111}{119.751}{4.736}{8.518}{r}{262848}{5}
  \Texto{99.152}{130.71}{66.747}{8.518}{r}{262848}{Associação linear}
  \Texto{192.91}{130.71}{85.265}{8.518}{r}{262848}{correlação de Pearson}
  \Texto{337.111}{130.71}{4.736}{8.518}{r}{262848}{6}
  \Texto{99.152}{141.669}{72.923}{8.518}{r}{262848}{Chance de eventos}
  \Texto{192.91}{141.669}{122.063}{8.518}{r}{262848}{probabilidade, binomial, Poisson}
  \Texto{337.111}{141.669}{4.736}{8.518}{r}{262848}{7}
  \Texto{49.227}{171.259}{357.083}{8.518}{r}{000000}{Medidas descritivas respondem a perguntas sobre valor típico, dispersão e associação. Mode-}
  \Texto{49.227}{182.218}{274.97}{8.518}{r}{000000}{los de probabilidade respondem a perguntas sobre a chance de eventos.}
  \Texto{37.21}{247.463}{40.499}{6.625}{b}{737696}{EMASC 2026}
  \Texto{80.653}{247.463}{2.319}{6.625}{r}{FFB014}{\textbullet{}}
  \Texto{85.913}{247.463}{117.821}{6.625}{r}{737696}{Liga Acadêmica Matemática \& Empresa}
  \Texto{206.682}{247.463}{2.319}{6.625}{r}{FFB014}{\textbullet{}}
  \Texto{211.943}{247.463}{13.25}{6.625}{r}{737696}{UFS}
  \Texto{397.549}{248.132}{7.367}{6.625}{b}{2C2084}{12}
  \Texto{406.02}{248.132}{10.315}{6.625}{r}{737696}{/ 40}
  \end{scope}
  \end{tikzpicture}
\end{frame}


\section{Análise descritiva}

% Página 17 do PDF de Gabriel
\begin{frame}[plain,noframenumbering,label=original-17]
  \begin{tikzpicture}[remember picture,overlay,x=1bp,y=-1bp]
  \begin{scope}[shift={(current page.north west)}]
  \Fundo{17}
  \Texto{44.456}{41.894}{90.657}{8.518}{b}{34E0C8}{ANÁLISE DESCRITIVA}
  \Texto{43.039}{116.125}{63.147}{56.787}{b}{FFB014}{04}
  \Texto{43.039}{155.148}{209.804}{24.608}{b}{FFFFFF}{Análise descritiva}
  \Texto{43.039}{234.662}{46.285}{7.572}{b}{34E0C8}{EMASC 2026}
  \Texto{91.429}{234.662}{2.65}{7.572}{b}{FFB014}{\textbullet{}}
  \Texto{96.176}{234.662}{21.45}{7.572}{b}{34E0C8}{LAME}
  \Link{42.042}{135.586}{211.795}{20.866}{original-17}
  \end{scope}
  \end{tikzpicture}
\end{frame}


% Página 18 do PDF de Gabriel
\begin{frame}[plain,noframenumbering,label=original-18]
  \begin{tikzpicture}[remember picture,overlay,x=1bp,y=-1bp]
  \begin{scope}[shift={(current page.north west)}]
  \Fundo{18}
  \Texto{40.175}{29.916}{253.823}{17.983}{b}{2C2084}{Medidas de tendência central}
  \Texto{35.433}{63.506}{295.975}{8.518}{i}{737696}{Uma lista de 15.000 valores não se comunica; um número representativo, sim.}
  \Texto{34.223}{100.879}{26.822}{9.465}{b}{CC8D10}{Média}
  \Texto{34.223}{115.823}{88.001}{9.465}{r}{000000}{O ponto de equilíbrio}
  \Texto{34.223}{127.778}{50.645}{9.465}{r}{000000}{dos valores.}
  \Texto{34.223}{145.711}{32.198}{11.357}{b}{000000}{4,57 h}
  \Texto{177.266}{100.642}{37.867}{9.465}{b}{027499}{Mediana}
  \Texto{177.266}{115.586}{89.089}{9.465}{r}{000000}{O valor central da fila}
  \Texto{177.266}{127.542}{42.619}{9.465}{r}{000000}{ordenada.}
  \Texto{177.266}{145.474}{16.411}{11.357}{b}{000000}{4 h}
  \Texto{320.309}{100.439}{24.712}{9.465}{b}{007078}{Moda}
  \Texto{320.309}{115.383}{83.922}{9.465}{r}{000000}{O valor que mais se}
  \Texto{320.309}{127.338}{29.321}{9.465}{r}{000000}{repete.}
  \Texto{320.309}{145.271}{16.411}{11.357}{b}{000000}{4 h}
  \Texto{49.227}{177.332}{196.246}{8.518}{r}{000000}{Horas diárias em redes sociais de sete estudantes:}
  \Texto{249.274}{177.332}{69.899}{8.518}{b}{000000}{2, 3, 4, 4, 5, 6 e 8}
  \Texto{319.172}{177.332}{87.139}{8.518}{r}{000000}{. Com poucos valores,}
  \Texto{49.227}{188.291}{357.083}{8.518}{r}{000000}{cada cálculo pode ser conferido manualmente; o procedimento é o mesmo para 15.000 regis-}
  \Texto{49.227}{199.25}{16.44}{8.518}{r}{000000}{tros.}
  \Texto{37.21}{247.463}{40.499}{6.625}{b}{737696}{EMASC 2026}
  \Texto{80.653}{247.463}{2.319}{6.625}{r}{FFB014}{\textbullet{}}
  \Texto{85.913}{247.463}{117.821}{6.625}{r}{737696}{Liga Acadêmica Matemática \& Empresa}
  \Texto{206.682}{247.463}{2.319}{6.625}{r}{FFB014}{\textbullet{}}
  \Texto{211.943}{247.463}{13.25}{6.625}{r}{737696}{UFS}
  \Texto{397.549}{248.132}{7.367}{6.625}{b}{2C2084}{13}
  \Texto{406.02}{248.132}{10.315}{6.625}{r}{737696}{/ 40}
  \end{scope}
  \end{tikzpicture}
\end{frame}


% Página 19 do PDF de Gabriel
\begin{frame}[plain,noframenumbering,label=original-19]
  \begin{tikzpicture}[remember picture,overlay,x=1bp,y=-1bp]
  \begin{scope}[shift={(current page.north west)}]
  \Fundo{19}
  \Texto{40.175}{28.518}{204.857}{17.983}{b}{2C2084}{Média, mediana e moda}
  \Texto{85.806}{83.374}{5.262}{9.465}{i}{262848}{n}
  \Texto{99.379}{64.087}{3.684}{6.625}{i}{262848}{n}
  \Texto{251.939}{87.561}{129.622}{6.974}{t}{000000}{x\ =\ np.array([2,\ 3,\ 4,\ 4,\ 5,\ 6,\ 8])}
  \Texto{251.939}{105.494}{29.625}{6.974}{t}{000000}{x.mean()}
  \Texto{337.13}{105.494}{29.625}{6.974}{c}{737696}{\#\ 4.5714}
  \Texto{251.939}{114.46}{44.437}{6.974}{t}{000000}{np.median(x)}
  \Texto{337.13}{114.46}{18.515}{6.974}{c}{737696}{\#\ 4.0}
  \Texto{251.939}{123.427}{70.359}{6.974}{t}{000000}{pd.Series(x).mode()}
  \Texto{337.13}{123.427}{11.109}{6.974}{c}{737696}{\#\ 4}
  \Texto{49.227}{168.935}{36.287}{8.518}{r}{000000}{regra que}
  \Texto{87.881}{168.935}{42.366}{8.966}{t}{2C2084}{np.median}
  \Texto{132.605}{168.935}{201.655}{8.518}{r}{000000}{aplica automaticamente. A moda pode não ser única.}
  \Texto{37.21}{247.463}{40.499}{6.625}{b}{737696}{EMASC 2026}
  \Texto{80.653}{247.463}{2.319}{6.625}{r}{FFB014}{\textbullet{}}
  \Texto{85.913}{247.463}{117.821}{6.625}{r}{737696}{Liga Acadêmica Matemática \& Empresa}
  \Texto{206.682}{247.463}{2.319}{6.625}{r}{FFB014}{\textbullet{}}
  \Texto{211.943}{247.463}{13.25}{6.625}{r}{737696}{UFS}
  \Texto{397.549}{248.132}{7.367}{6.625}{b}{2C2084}{14}
  \Texto{406.02}{248.132}{10.315}{6.625}{r}{737696}{/ 40}
  \end{scope}
  \end{tikzpicture}
\end{frame}


% Páginas 20 a 21 do PDF de Gabriel
\begin{frame}[plain,noframenumbering,label=original-20]
  \begin{tikzpicture}[remember picture,overlay,x=1bp,y=-1bp]
  \begin{scope}[shift={(current page.north west)}]
  \only<1>{\Fundo{20}}
  \only<2>{\Fundo{21}}
  \Texto{40.175}{29.916}{221.239}{17.983}{b}{2C2084}{Média e valores extremos}
  \Texto{35.433}{63.506}{295.754}{8.518}{i}{737696}{Uma medida só é confiável se soubermos como ela reage a valores extremos.}
  \only<1>{\Texto{49.227}{226.584}{7.104}{8.518}{r}{000000}{4.}}
  \Texto{37.21}{247.463}{40.499}{6.625}{b}{737696}{EMASC 2026}
  \Texto{80.653}{247.463}{2.319}{6.625}{r}{FFB014}{\textbullet{}}
  \Texto{85.913}{247.463}{117.821}{6.625}{r}{737696}{Liga Acadêmica Matemática \& Empresa}
  \Texto{206.682}{247.463}{2.319}{6.625}{r}{FFB014}{\textbullet{}}
  \Texto{211.943}{247.463}{13.25}{6.625}{r}{737696}{UFS}
  \Texto{397.549}{248.132}{7.367}{6.625}{b}{2C2084}{15}
  \Texto{406.02}{248.132}{10.315}{6.625}{r}{737696}{/ 40}
  \only<2>{\Texto{49.227}{226.584}{38.348}{8.518}{r}{000000}{A média é}}
  \only<2>{\Texto{89.942}{226.584}{33.621}{8.518}{b}{000000}{sensível}}
  \only<2>{\Texto{125.93}{226.584}{136.961}{8.518}{r}{000000}{a valores extremos; a mediana, não.}}
  \end{scope}
  \end{tikzpicture}
\end{frame}


% Página 22 do PDF de Gabriel
\begin{frame}[plain,noframenumbering,label=original-22]
  \begin{tikzpicture}[remember picture,overlay,x=1bp,y=-1bp]
  \begin{scope}[shift={(current page.north west)}]
  \Fundo{22}
  \Texto{40.175}{27.942}{124.906}{17.983}{b}{2C2084}{Desvio padrão}
  \Texto{35.433}{63.506}{264.723}{8.518}{i}{737696}{O valor típico não diz se os dados estão concentrados ou espalhados.}
  \Texto{49.227}{209.394}{357.083}{8.518}{r}{000000}{A soma dos desvios em relação à média é zero. Elevar ao quadrado elimina o sinal e a raiz}
  \Texto{37.21}{247.463}{40.499}{6.625}{b}{737696}{EMASC 2026}
  \Texto{80.653}{247.463}{2.319}{6.625}{r}{FFB014}{\textbullet{}}
  \Texto{85.913}{247.463}{117.821}{6.625}{r}{737696}{Liga Acadêmica Matemática \& Empresa}
  \Texto{206.682}{247.463}{2.319}{6.625}{r}{FFB014}{\textbullet{}}
  \Texto{211.943}{247.463}{13.25}{6.625}{r}{737696}{UFS}
  \Texto{397.549}{248.132}{7.367}{6.625}{b}{2C2084}{16}
  \Texto{406.02}{248.132}{10.315}{6.625}{r}{737696}{/ 40}
  \end{scope}
  \end{tikzpicture}
\end{frame}


% Página 23 do PDF de Gabriel
\begin{frame}[plain,noframenumbering,label=original-23]
  \begin{tikzpicture}[remember picture,overlay,x=1bp,y=-1bp]
  \begin{scope}[shift={(current page.north west)}]
  \Fundo{23}
  \Texto{40.175}{28.167}{223.792}{17.983}{b}{2C2084}{Variância e desvio padrão}
  \Texto{102.646}{69.8}{5.262}{9.465}{r}{262848}{1}
  \Texto{124.86}{64.087}{3.684}{6.625}{i}{262848}{n}
  \Texto{113.253}{116.688}{5.262}{9.465}{r}{262848}{6}
  \Texto{251.939}{87.561}{59.256}{6.974}{t}{000000}{d\ =\ x\ -\ x.mean()}
  \Texto{251.939}{96.528}{25.922}{6.974}{t}{000000}{(d**2).}
  \Texto{277.861}{96.528}{11.109}{6.974}{t}{2C2084}{sum}
  \Texto{288.97}{96.528}{7.406}{6.974}{t}{000000}{()}
  \Texto{318.61}{96.528}{29.625}{6.974}{c}{737696}{\#\ 23.714}
  \Texto{251.939}{105.494}{48.14}{6.974}{t}{000000}{x.var(ddof=1)}
  \Texto{318.61}{105.494}{25.922}{6.974}{c}{737696}{\#\ 3.952}
  \Texto{251.939}{114.46}{48.14}{6.974}{t}{000000}{x.std(ddof=1)}
  \Texto{318.61}{114.46}{25.922}{6.974}{c}{737696}{\#\ 1.988}
  \Texto{49.227}{153.918}{44.49}{8.518}{r}{000000}{Em NumPy,}
  \Texto{96.008}{153.918}{28.244}{8.966}{t}{2C2084}{ddof=0}
  \Texto{126.518}{153.918}{211.664}{8.518}{r}{000000}{calcula a variância populacional; em pandas, o padrão é}
  \Texto{340.456}{153.918}{28.244}{8.966}{t}{2C2084}{ddof=1}
  \Texto{368.694}{153.918}{37.616}{8.518}{r}{000000}{, a variân-}
  \Texto{49.227}{164.877}{319.297}{8.518}{r}{000000}{cia amostral. Para estimar a variância da população a partir de uma amostra, usa-se}
  \Texto{370.888}{164.877}{28.244}{8.966}{t}{2C2084}{ddof=1}
  \Texto{399.125}{164.877}{2.368}{8.518}{r}{000000}{.}
  \Texto{37.21}{247.463}{40.499}{6.625}{b}{737696}{EMASC 2026}
  \Texto{80.653}{247.463}{2.319}{6.625}{r}{FFB014}{\textbullet{}}
  \Texto{85.913}{247.463}{117.821}{6.625}{r}{737696}{Liga Acadêmica Matemática \& Empresa}
  \Texto{206.682}{247.463}{2.319}{6.625}{r}{FFB014}{\textbullet{}}
  \Texto{211.943}{247.463}{13.25}{6.625}{r}{737696}{UFS}
  \Texto{397.549}{248.132}{7.367}{6.625}{b}{2C2084}{17}
  \Texto{406.02}{248.132}{10.315}{6.625}{r}{737696}{/ 40}
  \end{scope}
  \end{tikzpicture}
\end{frame}


% Página 24 do PDF de Gabriel
\begin{frame}[plain,noframenumbering,label=original-24]
  \begin{tikzpicture}[remember picture,overlay,x=1bp,y=-1bp]
  \begin{scope}[shift={(current page.north west)}]
  \Fundo{24}
  \Texto{35.433}{63.506}{282.891}{8.518}{i}{737696}{Estimar a variância da população a partir de uma amostra exige um ajuste.}
  \Texto{49.227}{181.15}{296.29}{8.518}{r}{000000}{A média é calculada da própria amostra, o que reduz os desvios. Dividir por}
  \Texto{348.536}{181.15}{4.736}{8.518}{i}{000000}{n}
  \Texto{356.428}{181.15}{42.13}{8.518}{b}{000000}{subestima}
  \Texto{401.573}{181.15}{4.736}{8.518}{r}{000000}{a}
  \Texto{37.21}{247.463}{40.499}{6.625}{b}{737696}{EMASC 2026}
  \Texto{80.653}{247.463}{2.319}{6.625}{r}{FFB014}{\textbullet{}}
  \Texto{85.913}{247.463}{117.821}{6.625}{r}{737696}{Liga Acadêmica Matemática \& Empresa}
  \Texto{206.682}{247.463}{2.319}{6.625}{r}{FFB014}{\textbullet{}}
  \Texto{211.943}{247.463}{13.25}{6.625}{r}{737696}{UFS}
  \Texto{397.549}{248.132}{7.367}{6.625}{b}{2C2084}{18}
  \Texto{406.02}{248.132}{10.315}{6.625}{r}{737696}{/ 40}
  \end{scope}
  \end{tikzpicture}
\end{frame}


% Página 25 do PDF de Gabriel
\begin{frame}[plain,noframenumbering,label=original-25]
  \begin{tikzpicture}[remember picture,overlay,x=1bp,y=-1bp]
  \begin{scope}[shift={(current page.north west)}]
  \Fundo{25}
  \Texto{40.175}{29.336}{117.282}{17.983}{b}{2C2084}{Quartis e IQR}
  \Texto{35.433}{63.506}{298.139}{8.518}{i}{737696}{Quando há valores extremos, é útil medir a dispersão sem depender da média.}
  \Texto{49.227}{142.749}{130.479}{8.518}{r}{000000}{respondem a 25\%, 50\% e 75\%. O}
  \Texto{182.209}{142.749}{15.145}{8.518}{b}{000000}{IQR}
  \Texto{199.859}{142.749}{138.179}{8.518}{r}{000000}{é a amplitude dos 50\% centrais e as}
  \Texto{340.541}{142.749}{26.823}{8.518}{b}{000000}{cercas}
  \Texto{369.87}{142.749}{36.44}{8.518}{r}{000000}{delimitam}
  \Texto{37.21}{247.463}{40.499}{6.625}{b}{737696}{EMASC 2026}
  \Texto{80.653}{247.463}{2.319}{6.625}{r}{FFB014}{\textbullet{}}
  \Texto{85.913}{247.463}{117.821}{6.625}{r}{737696}{Liga Acadêmica Matemática \& Empresa}
  \Texto{206.682}{247.463}{2.319}{6.625}{r}{FFB014}{\textbullet{}}
  \Texto{211.943}{247.463}{13.25}{6.625}{r}{737696}{UFS}
  \Texto{397.549}{248.132}{7.367}{6.625}{b}{2C2084}{19}
  \Texto{406.02}{248.132}{10.315}{6.625}{r}{737696}{/ 40}
  \end{scope}
  \end{tikzpicture}
\end{frame}


\section{Distribuições}

% Página 26 do PDF de Gabriel
\begin{frame}[plain,noframenumbering,label=original-26]
  \begin{tikzpicture}[remember picture,overlay,x=1bp,y=-1bp]
  \begin{scope}[shift={(current page.north west)}]
  \Fundo{26}
  \Texto{44.456}{41.894}{90.657}{8.518}{b}{34E0C8}{ANÁLISE DESCRITIVA}
  \Texto{43.039}{116.125}{63.147}{56.787}{b}{FFB014}{05}
  \Texto{43.039}{154.533}{155.397}{24.608}{b}{FFFFFF}{Distribuições}
  \Texto{43.039}{234.662}{46.285}{7.572}{b}{34E0C8}{EMASC 2026}
  \Texto{91.429}{234.662}{2.65}{7.572}{b}{FFB014}{\textbullet{}}
  \Texto{96.176}{234.662}{21.45}{7.572}{b}{34E0C8}{LAME}
  \Link{42.042}{135.586}{157.388}{25.418}{original-26}
  \end{scope}
  \end{tikzpicture}
\end{frame}


% Páginas 27 a 31 do PDF de Gabriel
\begin{frame}[plain,noframenumbering,label=original-27]
  \begin{tikzpicture}[remember picture,overlay,x=1bp,y=-1bp]
  \begin{scope}[shift={(current page.north west)}]
  \only<1>{\Fundo{27}}
  \only<2>{\Fundo{28}}
  \only<3>{\Fundo{29}}
  \only<4>{\Fundo{30}}
  \only<5>{\Fundo{31}}
  \Texto{40.175}{27.803}{198.275}{17.983}{b}{2C2084}{Construção do boxplot}
  \Texto{35.433}{63.506}{246.971}{8.518}{i}{737696}{Os quartis permitem um gráfico que resume a distribuição inteira.}
  \only<1>{\Texto{49.227}{215.166}{177.703}{8.518}{r}{000000}{Os mesmos sete valores, agora sobre um eixo.}}
  \Texto{37.21}{247.463}{40.499}{6.625}{b}{737696}{EMASC 2026}
  \Texto{80.653}{247.463}{2.319}{6.625}{r}{FFB014}{\textbullet{}}
  \Texto{85.913}{247.463}{117.821}{6.625}{r}{737696}{Liga Acadêmica Matemática \& Empresa}
  \Texto{206.682}{247.463}{2.319}{6.625}{r}{FFB014}{\textbullet{}}
  \Texto{211.943}{247.463}{13.25}{6.625}{r}{737696}{UFS}
  \Texto{397.549}{248.132}{7.367}{6.625}{b}{2C2084}{20}
  \Texto{406.02}{248.132}{10.315}{6.625}{r}{737696}{/ 40}
  \only<2>{\Texto{49.227}{215.166}{135.147}{8.518}{r}{000000}{Quartis: 25\% dos valores abaixo de}}
  \only<2>{\Texto{186.739}{215.166}{6.627}{8.518}{i}{000000}{Q}}
  \only<2>{\Texto{193.366}{216.162}{3.157}{5.679}{r}{000000}{1}}
  \only<2>{\Texto{197.022}{215.166}{161.936}{8.518}{r}{000000}{, 50\% abaixo da mediana e 75\% abaixo de}}
  \only<2>{\Texto{361.323}{215.166}{6.627}{8.518}{i}{000000}{Q}}
  \only<2>{\Texto{367.95}{216.162}{3.157}{5.679}{r}{000000}{3}}
  \only<2>{\Texto{371.606}{215.166}{2.368}{8.518}{r}{000000}{.}}
  \only<3>{\Texto{49.227}{215.166}{52.812}{8.518}{r}{000000}{A caixa vai de}}
  \only<3>{\Texto{104.406}{215.166}{6.627}{8.518}{i}{000000}{Q}}
  \only<3>{\Texto{111.033}{216.162}{3.157}{5.679}{r}{000000}{1}}
  \only<3>{\Texto{117.056}{215.166}{4.736}{8.518}{r}{000000}{a}}
  \only<3>{\Texto{124.16}{215.166}{6.627}{8.518}{i}{000000}{Q}}
  \only<3>{\Texto{130.787}{216.162}{3.157}{5.679}{r}{000000}{3}}
  \only<3>{\Texto{134.442}{215.166}{174.875}{8.518}{r}{000000}{; sua largura é o IQR e cobre os 50\% centrais.}}
  \only<4>{\Texto{49.227}{215.166}{182.728}{8.518}{r}{000000}{Os bigodes terminam no menor e no maior valor}}
  \only<4>{\Texto{234.321}{215.166}{26.329}{8.518}{b}{000000}{dentro}}
  \only<4>{\Texto{263.018}{215.166}{88.63}{8.518}{r}{000000}{das cercas: aqui, 2 e 8.}}
  \only<5>{\Texto{49.227}{215.327}{324.851}{8.518}{r}{000000}{Com 24 h no lugar de 8 h, o valor ultrapassa a cerca superior (8,5) e é marcado como}}
  \only<5>{\Texto{376.488}{215.327}{27.454}{8.518}{b}{000000}{atípico}}
  \only<5>{\Texto{403.941}{215.327}{2.368}{8.518}{r}{000000}{.}}
  \only<5>{\Texto{49.227}{226.286}{160.087}{8.518}{r}{000000}{A caixa e o IQR permanecem os mesmos.}}
  \end{scope}
  \end{tikzpicture}
\end{frame}


% Páginas 32 a 34 do PDF de Gabriel
\begin{frame}[plain,noframenumbering,label=original-32]
  \begin{tikzpicture}[remember picture,overlay,x=1bp,y=-1bp]
  \begin{scope}[shift={(current page.north west)}]
  \only<1>{\Fundo{32}}
  \only<2>{\Fundo{33}}
  \only<3>{\Fundo{34}}
  \Texto{40.175}{27.803}{230.787}{17.983}{b}{2C2084}{Construção do histograma}
  \Texto{35.433}{63.506}{304.297}{8.518}{i}{737696}{Para ver a forma da distribuição, contamos quantos valores caem em cada faixa.}
  \only<1>{\Texto{49.227}{209.496}{109.15}{8.518}{r}{000000}{Os sete valores sobre a reta.}}
  \only<1,2>{\Texto{37.21}{247.463}{40.499}{6.625}{b}{737696}{EMASC 2026}}
  \only<1,2>{\Texto{80.653}{247.463}{2.319}{6.625}{r}{FFB014}{\textbullet{}}}
  \only<1,2>{\Texto{85.913}{247.463}{117.821}{6.625}{r}{737696}{Liga Acadêmica Matemática \& Empresa}}
  \only<1,2>{\Texto{206.682}{247.463}{2.319}{6.625}{r}{FFB014}{\textbullet{}}}
  \only<1,2>{\Texto{211.943}{247.463}{13.25}{6.625}{r}{737696}{UFS}}
  \only<1,2>{\Texto{397.549}{248.132}{7.367}{6.625}{b}{2C2084}{21}}
  \only<1,2>{\Texto{406.02}{248.132}{10.315}{6.625}{r}{737696}{/ 40}}
  \only<2>{\Texto{49.227}{209.3}{357.083}{8.518}{r}{000000}{Faixas de mesma largura cobrem a amplitude dos dados; conta-se quantos valores caem em}}
  \only<2>{\Texto{49.227}{220.259}{40.938}{8.518}{r}{000000}{cada faixa.}}
  \only<3>{\Texto{49.227}{209.495}{357.083}{8.518}{r}{000000}{Uma barra por faixa, com altura igual à contagem. A largura da faixa é a amplitude dividida pelo}}
  \only<3>{\Texto{49.227}{220.454}{298.871}{8.518}{r}{000000}{número de faixas: faixas largas escondem o formato; estreitas produzem ruído.}}
  \only<3>{\Texto{37.21}{247.462}{40.499}{6.625}{b}{737696}{EMASC 2026}}
  \only<3>{\Texto{80.653}{247.462}{2.319}{6.625}{r}{FFB014}{\textbullet{}}}
  \only<3>{\Texto{85.913}{247.462}{117.821}{6.625}{r}{737696}{Liga Acadêmica Matemática \& Empresa}}
  \only<3>{\Texto{206.682}{247.462}{2.319}{6.625}{r}{FFB014}{\textbullet{}}}
  \only<3>{\Texto{211.943}{247.462}{13.25}{6.625}{r}{737696}{UFS}}
  \only<3>{\Texto{397.549}{248.131}{7.367}{6.625}{b}{2C2084}{21}}
  \only<3>{\Texto{406.02}{248.131}{10.315}{6.625}{r}{737696}{/ 40}}
  \end{scope}
  \end{tikzpicture}
\end{frame}


% Página 35 do PDF de Gabriel
\begin{frame}[plain,noframenumbering,label=original-35]
  \begin{tikzpicture}[remember picture,overlay,x=1bp,y=-1bp]
  \begin{scope}[shift={(current page.north west)}]
  \Fundo{35}
  \Texto{40.175}{27.803}{219.315}{17.983}{b}{2C2084}{Distribuições observadas}
  \Texto{35.433}{63.506}{280.285}{8.518}{i}{737696}{A forma da distribuição explica quando a média representa bem os dados.}
  \Texto{49.227}{203.827}{357.083}{8.518}{r}{000000}{Em distribuições assimétricas, a média é deslocada no sentido da cauda: na pontuação de}
  \Texto{49.227}{225.745}{225.616}{8.518}{r}{000000}{sono e de redes, média e mediana praticamente coincidem.}
  \Texto{37.21}{247.463}{40.499}{6.625}{b}{737696}{EMASC 2026}
  \Texto{80.653}{247.463}{2.319}{6.625}{r}{FFB014}{\textbullet{}}
  \Texto{85.913}{247.463}{117.821}{6.625}{r}{737696}{Liga Acadêmica Matemática \& Empresa}
  \Texto{206.682}{247.463}{2.319}{6.625}{r}{FFB014}{\textbullet{}}
  \Texto{211.943}{247.463}{13.25}{6.625}{r}{737696}{UFS}
  \Texto{397.549}{248.132}{7.367}{6.625}{b}{2C2084}{22}
  \Texto{406.02}{248.132}{10.315}{6.625}{r}{737696}{/ 40}
  \end{scope}
  \end{tikzpicture}
\end{frame}


% Página 36 do PDF de Gabriel
\begin{frame}[plain,noframenumbering,label=original-36]
  \begin{tikzpicture}[remember picture,overlay,x=1bp,y=-1bp]
  \begin{scope}[shift={(current page.north west)}]
  \Fundo{36}
  \Texto{40.175}{28.028}{215.844}{17.983}{b}{2C2084}{Comparação entre níveis}
  \Texto{35.433}{63.506}{225.276}{8.518}{i}{737696}{As mesmas medidas e gráficos permitem comparar grupos.}
  \Texto{49.227}{203.827}{357.083}{8.518}{r}{000000}{As médias de uso de IA diferem em centésimos de hora, enquanto o desvio padrão é de cerca}
  \Texto{49.227}{214.786}{86.296}{8.518}{r}{000000}{de 1,7 h: a diferença é}
  \Texto{138.092}{214.786}{117.881}{8.518}{b}{000000}{pequena diante da dispersão}
  \Texto{255.97}{214.786}{150.343}{8.518}{r}{000000}{. Além disso, os níveis correspondem a}
  \Texto{49.227}{225.745}{85.904}{8.518}{r}{000000}{faixas etárias distintas.}
  \Texto{37.21}{247.463}{40.499}{6.625}{b}{737696}{EMASC 2026}
  \Texto{80.653}{247.463}{2.319}{6.625}{r}{FFB014}{\textbullet{}}
  \Texto{85.913}{247.463}{117.821}{6.625}{r}{737696}{Liga Acadêmica Matemática \& Empresa}
  \Texto{206.682}{247.463}{2.319}{6.625}{r}{FFB014}{\textbullet{}}
  \Texto{211.943}{247.463}{13.25}{6.625}{r}{737696}{UFS}
  \Texto{397.549}{248.132}{7.367}{6.625}{b}{2C2084}{23}
  \Texto{406.02}{248.132}{10.315}{6.625}{r}{737696}{/ 40}
  \end{scope}
  \end{tikzpicture}
\end{frame}


\section{Correlação linear}

% Página 37 do PDF de Gabriel
\begin{frame}[plain,noframenumbering,label=original-37]
  \begin{tikzpicture}[remember picture,overlay,x=1bp,y=-1bp]
  \begin{scope}[shift={(current page.north west)}]
  \Fundo{37}
  \Texto{44.456}{41.894}{90.657}{8.518}{b}{34E0C8}{ANÁLISE DESCRITIVA}
  \Texto{43.039}{116.125}{63.147}{56.787}{b}{FFB014}{06}
  \Texto{43.039}{154.533}{202.078}{24.608}{b}{FFFFFF}{Correlação linear}
  \Texto{43.039}{234.662}{46.285}{7.572}{b}{34E0C8}{EMASC 2026}
  \Texto{91.429}{234.662}{2.65}{7.572}{b}{FFB014}{\textbullet{}}
  \Texto{96.176}{234.662}{21.45}{7.572}{b}{34E0C8}{LAME}
  \Link{42.042}{135.586}{204.068}{25.418}{original-37}
  \end{scope}
  \end{tikzpicture}
\end{frame}


% Página 38 do PDF de Gabriel
\begin{frame}[plain,noframenumbering,label=original-38]
  \begin{tikzpicture}[remember picture,overlay,x=1bp,y=-1bp]
  \begin{scope}[shift={(current page.north west)}]
  \Fundo{38}
  \Texto{40.175}{27.803}{236.128}{17.983}{b}{2C2084}{Por que medir associação?}
  \Texto{35.433}{63.506}{307.909}{8.518}{i}{737696}{Duas variáveis podem variar juntas; queremos resumir isso em um único número.}
  \Texto{38.534}{100.643}{34.716}{9.465}{b}{027499}{Direção}
  \Texto{38.534}{115.586}{163.708}{9.465}{r}{000000}{Quando uma variável aumenta, a outra}
  \Texto{38.534}{127.542}{136.156}{9.465}{r}{000000}{tende a aumentar ou a diminuir?}
  \Texto{250.943}{100.643}{52.594}{9.465}{b}{CC8D10}{Intensidade}
  \Texto{250.943}{115.587}{143.084}{9.465}{r}{000000}{Quão próxima de uma reta é essa}
  \Texto{250.943}{127.542}{36.296}{9.465}{r}{000000}{relação?}
  \Texto{49.227}{162.195}{153.528}{8.518}{r}{000000}{O coeficiente de correlação de Pearson,}
  \Texto{205.623}{162.195}{2.836}{8.518}{i}{000000}{r}
  \Texto{209.418}{162.195}{196.894}{8.518}{r}{000000}{, resume esses dois aspectos em um número entre}
  \Texto{49.227}{184.113}{55.87}{8.518}{r}{000000}{saúde mental?}
  \Texto{37.21}{247.463}{40.499}{6.625}{b}{737696}{EMASC 2026}
  \Texto{80.653}{247.463}{2.319}{6.625}{r}{FFB014}{\textbullet{}}
  \Texto{85.913}{247.463}{117.821}{6.625}{r}{737696}{Liga Acadêmica Matemática \& Empresa}
  \Texto{206.682}{247.463}{2.319}{6.625}{r}{FFB014}{\textbullet{}}
  \Texto{211.943}{247.463}{13.25}{6.625}{r}{737696}{UFS}
  \Texto{397.549}{248.132}{7.367}{6.625}{b}{2C2084}{24}
  \Texto{406.02}{248.132}{10.315}{6.625}{r}{737696}{/ 40}
  \end{scope}
  \end{tikzpicture}
\end{frame}


% Página 39 do PDF de Gabriel
\begin{frame}[plain,noframenumbering,label=original-39]
  \begin{tikzpicture}[remember picture,overlay,x=1bp,y=-1bp]
  \begin{scope}[shift={(current page.north west)}]
  \Fundo{39}
  \Texto{40.175}{27.803}{197.052}{17.983}{b}{2C2084}{Correlação de Pearson}
  \Texto{35.433}{63.506}{225.523}{8.518}{i}{737696}{A fórmula traduz a associação em um número padronizado.}
  \Texto{46.238}{128.971}{107.397}{6.974}{t}{000000}{a\ =\ np.array([1,\ 2,\ 3,\ 4,\ 5])}
  \Texto{46.238}{137.938}{107.397}{6.974}{t}{000000}{b\ =\ np.array([2,\ 1,\ 4,\ 3,\ 5])}
  \Texto{46.238}{146.904}{129.622}{6.974}{t}{000000}{da,\ db\ =\ a\ -\ a.mean(),\ b\ -\ b.mean()}
  \Texto{46.238}{155.87}{51.843}{6.974}{t}{000000}{num\ =\ (da*db).}
  \Texto{98.081}{155.87}{11.116}{6.974}{t}{2C2084}{sum}
  \Texto{109.197}{155.87}{7.406}{6.974}{t}{000000}{()}
  \Texto{212.915}{155.87}{11.109}{6.974}{c}{737696}{\#\ 8}
  \Texto{46.238}{164.837}{81.475}{6.974}{t}{000000}{den\ =\ np.sqrt((da**2).}
  \Texto{127.713}{164.837}{11.109}{6.974}{t}{2C2084}{sum}
  \Texto{138.822}{164.837}{40.734}{6.974}{t}{000000}{()*(db**2).}
  \Texto{179.556}{164.837}{11.109}{6.974}{t}{2C2084}{sum}
  \Texto{190.665}{164.837}{11.109}{6.974}{t}{000000}{())}
  \Texto{212.915}{164.837}{14.812}{6.974}{c}{737696}{\#\ 10}
  \Texto{46.238}{173.803}{33.328}{6.974}{t}{000000}{num\ /\ den}
  \Texto{212.915}{173.803}{18.515}{6.974}{c}{737696}{\#\ 0.8}
  \Texto{46.238}{182.769}{85.178}{6.974}{t}{000000}{np.corrcoef(a,\ b)[0,\ 1]}
  \Texto{212.915}{182.769}{18.515}{6.974}{c}{737696}{\#\ 0.8}
  \Texto{49.227}{217.319}{350.78}{8.518}{r}{000000}{O numerador soma os produtos dos desvios; o denominador padroniza o resultado. Por isso}
  \Texto{402.515}{217.319}{2.836}{8.518}{i}{000000}{r}
  \Texto{37.21}{247.463}{40.499}{6.625}{b}{737696}{EMASC 2026}
  \Texto{80.653}{247.463}{2.319}{6.625}{r}{FFB014}{\textbullet{}}
  \Texto{85.913}{247.463}{117.821}{6.625}{r}{737696}{Liga Acadêmica Matemática \& Empresa}
  \Texto{206.682}{247.463}{2.319}{6.625}{r}{FFB014}{\textbullet{}}
  \Texto{211.943}{247.463}{13.25}{6.625}{r}{737696}{UFS}
  \Texto{397.549}{248.132}{7.367}{6.625}{b}{2C2084}{25}
  \Texto{406.02}{248.132}{10.315}{6.625}{r}{737696}{/ 40}
  \end{scope}
  \end{tikzpicture}
\end{frame}


% Páginas 40 a 43 do PDF de Gabriel
\begin{frame}[plain,noframenumbering,label=original-40]
  \begin{tikzpicture}[remember picture,overlay,x=1bp,y=-1bp]
  \begin{scope}[shift={(current page.north west)}]
  \only<1>{\Fundo{40}}
  \only<2>{\Fundo{41}}
  \only<3>{\Fundo{42}}
  \only<4>{\Fundo{43}}
  \Texto{40.175}{27.803}{221.832}{17.983}{b}{2C2084}{Limitações do coeficiente}
  \Texto{35.433}{63.506}{240.361}{8.518}{i}{737696}{r resume associação linear; fora desse contexto, pode enganar.}
  \only<1>{\Texto{49.227}{207.561}{22.036}{8.518}{r}{000000}{sável.}}
  \Texto{37.21}{247.463}{40.499}{6.625}{b}{737696}{EMASC 2026}
  \Texto{80.653}{247.463}{2.319}{6.625}{r}{FFB014}{\textbullet{}}
  \Texto{85.913}{247.463}{117.821}{6.625}{r}{737696}{Liga Acadêmica Matemática \& Empresa}
  \Texto{206.682}{247.463}{2.319}{6.625}{r}{FFB014}{\textbullet{}}
  \Texto{211.943}{247.463}{13.25}{6.625}{r}{737696}{UFS}
  \Texto{397.549}{248.132}{7.367}{6.625}{b}{2C2084}{26}
  \Texto{406.02}{248.132}{10.315}{6.625}{r}{737696}{/ 40}
  \only<2,3>{\Texto{49.227}{196.603}{7.104}{8.518}{b}{000000}{2.}}
  \only<2>{\Texto{59.261}{196.603}{236.443}{8.518}{r}{000000}{Sem o ponto destacado, não há associação entre as variáveis.}}
  \only<3>{\Texto{59.261}{196.603}{191.553}{8.518}{r}{000000}{Um único ponto extremo cria uma correlação forte.}}
  \only<4>{\Texto{49.227}{196.603}{7.104}{8.518}{b}{000000}{3.}}
  \only<4>{\Texto{59.185}{196.603}{347.126}{8.518}{r}{000000}{O calor aumenta as vendas de sorvete e o número de afogamentos. A associação entre eles}}
  \only<4>{\Texto{49.227}{207.561}{50.665}{8.518}{r}{000000}{não é causal.}}
  \end{scope}
  \end{tikzpicture}
\end{frame}


% Página 44 do PDF de Gabriel
\begin{frame}[plain,noframenumbering,label=original-44]
  \begin{tikzpicture}[remember picture,overlay,x=1bp,y=-1bp]
  \begin{scope}[shift={(current page.north west)}]
  \Fundo{44}
  \Texto{40.175}{27.803}{190.866}{17.983}{b}{2C2084}{Correlação nos dados}
  \Texto{35.433}{63.506}{221.689}{8.518}{i}{737696}{Retomamos a pergunta inicial com o coeficiente calculado.}
  \Texto{49.227}{214.786}{357.083}{8.518}{r}{000000}{negativa de intensidade moderada. O coeficiente não indica quantos pontos se perdem por}
  \Texto{49.227}{225.745}{186.817}{8.518}{r}{000000}{hora de uso; mede apenas direção e intensidade.}
  \Texto{37.21}{247.463}{40.499}{6.625}{b}{737696}{EMASC 2026}
  \Texto{80.653}{247.463}{2.319}{6.625}{r}{FFB014}{\textbullet{}}
  \Texto{85.913}{247.463}{117.821}{6.625}{r}{737696}{Liga Acadêmica Matemática \& Empresa}
  \Texto{206.682}{247.463}{2.319}{6.625}{r}{FFB014}{\textbullet{}}
  \Texto{211.943}{247.463}{13.25}{6.625}{r}{737696}{UFS}
  \Texto{397.549}{248.132}{7.367}{6.625}{b}{2C2084}{27}
  \Texto{406.02}{248.132}{10.315}{6.625}{r}{737696}{/ 40}
  \end{scope}
  \end{tikzpicture}
\end{frame}


% Páginas 45 a 47 do PDF de Gabriel
\begin{frame}[plain,noframenumbering,label=original-45]
  \begin{tikzpicture}[remember picture,overlay,x=1bp,y=-1bp]
  \begin{scope}[shift={(current page.north west)}]
  \only<1>{\Fundo{45}}
  \only<2>{\Fundo{46}}
  \only<3>{\Fundo{47}}
  \Texto{40.175}{27.803}{185.867}{17.983}{b}{2C2084}{Matriz de correlações}
  \Texto{35.433}{63.506}{281.179}{8.518}{i}{737696}{O mesmo coeficiente pode ser calculado para todos os pares de variáveis.}
  \Texto{254.928}{93.035}{17.982}{8.518}{b}{132482}{Azul}
  \Texto{272.909}{93.035}{2.368}{8.518}{r}{000000}{:}
  \Texto{282.143}{93.035}{81.287}{8.518}{r}{000000}{associação negativa.}
  \Texto{372.259}{93.035}{38.391}{8.518}{b}{E27800}{Vermelho}
  \Texto{410.649}{93.035}{2.368}{8.518}{r}{000000}{:}
  \Texto{254.928}{103.994}{76.347}{8.518}{r}{000000}{associação positiva.}
  \Texto{335.224}{103.994}{29.345}{8.518}{b}{000000}{Branco}
  \Texto{364.569}{103.994}{48.45}{8.518}{r}{000000}{: próximo de}
  \Texto{254.928}{114.953}{18.467}{8.518}{r}{000000}{zero.}
  \Texto{37.209}{247.463}{40.499}{6.625}{b}{737696}{EMASC 2026}
  \Texto{80.653}{247.463}{2.319}{6.625}{r}{FFB014}{\textbullet{}}
  \Texto{85.914}{247.463}{117.821}{6.625}{r}{737696}{Liga Acadêmica Matemática \& Empresa}
  \Texto{206.683}{247.463}{2.319}{6.625}{r}{FFB014}{\textbullet{}}
  \Texto{211.943}{247.463}{13.25}{6.625}{r}{737696}{UFS}
  \Texto{397.549}{248.132}{7.367}{6.625}{b}{2C2084}{28}
  \Texto{406.02}{248.132}{10.315}{6.625}{r}{737696}{/ 40}
  \only<2,3>{\Texto{254.928}{140.933}{158.086}{8.518}{r}{000000}{A diagonal vale 1, pois cada variável é per-}}
  \only<2,3>{\Texto{254.928}{151.892}{158.086}{8.518}{r}{000000}{feitamente associada a si mesma, e a ma-}}
  \only<2,3>{\Texto{254.928}{162.851}{57.522}{8.518}{r}{000000}{triz é simétrica.}}
  \only<3>{\Texto{254.928}{188.635}{158.086}{8.518}{r}{000000}{A idade não apresenta associação com as}}
  \only<3>{\Texto{254.928}{199.594}{158.086}{8.518}{r}{000000}{demais variáveis, o que reforça a hipótese}}
  \only<3>{\Texto{254.928}{210.553}{77.906}{8.518}{r}{000000}{de dados simulados.}}
  \end{scope}
  \end{tikzpicture}
\end{frame}


% Páginas 48 a 51 do PDF de Gabriel
\begin{frame}[plain,noframenumbering,label=original-48]
  \begin{tikzpicture}[remember picture,overlay,x=1bp,y=-1bp]
  \begin{scope}[shift={(current page.north west)}]
  \only<1>{\Fundo{48}}
  \only<2>{\Fundo{49}}
  \only<3>{\Fundo{50}}
  \only<4>{\Fundo{51}}
  \Texto{40.175}{27.803}{224.871}{17.983}{b}{2C2084}{Associação e causalidade}
  \Texto{35.433}{63.506}{203.12}{8.518}{i}{737696}{Medir associação não responde à pergunta de causa.}
  \Texto{71.408}{84.712}{22.726}{8.518}{b}{262848}{redes}
  \Texto{55.552}{127.429}{54.439}{8.518}{b}{262848}{saúde mental}
  \Texto{35.302}{147.661}{94.938}{9.465}{r}{262848}{Redes afetam a saúde}
  \Texto{65.939}{159.616}{33.665}{9.465}{r}{262848}{mental?}
  \Texto{37.21}{247.464}{40.499}{6.625}{b}{737696}{EMASC 2026}
  \Texto{80.653}{247.464}{2.319}{6.625}{r}{FFB014}{\textbullet{}}
  \Texto{85.913}{247.464}{117.821}{6.625}{r}{737696}{Liga Acadêmica Matemática \& Empresa}
  \Texto{206.682}{247.464}{2.319}{6.625}{r}{FFB014}{\textbullet{}}
  \Texto{211.942}{247.464}{13.25}{6.625}{r}{737696}{UFS}
  \Texto{397.549}{248.133}{7.367}{6.625}{b}{2C2084}{29}
  \Texto{406.02}{248.133}{10.315}{6.625}{r}{737696}{/ 40}
  \only<2,3,4>{\Texto{215.409}{84.713}{22.726}{8.518}{b}{262848}{redes}}
  \only<2,3,4>{\Texto{199.553}{127.429}{54.439}{8.518}{b}{262848}{saúde mental}}
  \only<2,3,4>{\Texto{167.899}{147.661}{117.748}{9.465}{r}{262848}{A saúde mental muda o uso}}
  \only<2,3,4>{\Texto{205.728}{159.616}{42.089}{9.465}{r}{262848}{de redes?}}
  \only<3,4>{\Texto{355.394}{84.896}{26.977}{8.518}{b}{262848}{fator X}}
  \only<3,4>{\Texto{332.008}{132.993}{22.726}{8.518}{b}{262848}{redes}}
  \only<3,4>{\Texto{384.92}{133.191}{24.617}{8.518}{b}{262848}{saúde}}
  \only<3,4>{\Texto{317.668}{153.422}{106.211}{9.465}{r}{262848}{Um terceiro fator afeta as}}
  \only<3,4>{\Texto{357.882}{165.377}{25.781}{9.465}{r}{262848}{duas?}}
  \only<4>{\Texto{49.227}{185.544}{118.289}{8.518}{r}{000000}{Os três diagramas produzem o}}
  \only<4>{\Texto{170.3}{185.544}{29.822}{8.518}{b}{000000}{mesmo}}
  \only<4>{\Texto{202.904}{185.544}{203.402}{8.518}{r}{000000}{coeficiente de correlação. Dados observacionais, por}}
  \only<4>{\Texto{49.227}{196.503}{152.992}{8.518}{r}{000000}{si só, não permitem distinguir entre eles.}}
  \end{scope}
  \end{tikzpicture}
\end{frame}


\section{Probabilidade}

% Página 52 do PDF de Gabriel
\begin{frame}[plain,noframenumbering,label=original-52]
  \begin{tikzpicture}[remember picture,overlay,x=1bp,y=-1bp]
  \begin{scope}[shift={(current page.north west)}]
  \Fundo{52}
  \Texto{44.456}{41.894}{123.698}{8.518}{b}{34E0C8}{PROBABILIDADE E MODELOS}
  \Texto{43.039}{116.125}{63.147}{56.787}{b}{FFB014}{07}
  \Texto{43.039}{154.533}{162.238}{24.608}{b}{FFFFFF}{Probabilidade}
  \Texto{43.039}{234.662}{46.285}{7.572}{b}{34E0C8}{EMASC 2026}
  \Texto{91.429}{234.662}{2.65}{7.572}{b}{FFB014}{\textbullet{}}
  \Texto{96.176}{234.662}{21.45}{7.572}{b}{34E0C8}{LAME}
  \Link{42.042}{135.586}{164.229}{20.251}{original-52}
  \end{scope}
  \end{tikzpicture}
\end{frame}


% Página 53 do PDF de Gabriel
\begin{frame}[plain,noframenumbering,label=original-53]
  \begin{tikzpicture}[remember picture,overlay,x=1bp,y=-1bp]
  \begin{scope}[shift={(current page.north west)}]
  \Fundo{53}
  \Texto{40.175}{27.942}{178.584}{17.983}{b}{2C2084}{Conjuntos e eventos}
  \Texto{35.433}{63.506}{256.256}{8.518}{i}{737696}{Probabilidade fala de eventos; eventos são conjuntos de resultados.}
  \Texto{73.648}{116.403}{5.682}{8.518}{i}{262848}{A}
  \Texto{138.655}{116.403}{5.682}{8.518}{i}{262848}{B}
  \Texto{37.21}{247.463}{40.499}{6.625}{b}{737696}{EMASC 2026}
  \Texto{80.653}{247.463}{2.319}{6.625}{r}{FFB014}{\textbullet{}}
  \Texto{85.913}{247.463}{117.821}{6.625}{r}{737696}{Liga Acadêmica Matemática \& Empresa}
  \Texto{206.682}{247.463}{2.319}{6.625}{r}{FFB014}{\textbullet{}}
  \Texto{211.943}{247.463}{13.25}{6.625}{r}{737696}{UFS}
  \Texto{397.549}{248.132}{7.367}{6.625}{b}{2C2084}{30}
  \Texto{406.02}{248.132}{10.315}{6.625}{r}{737696}{/ 40}
  \end{scope}
  \end{tikzpicture}
\end{frame}


% Página 54 do PDF de Gabriel
\begin{frame}[plain,noframenumbering,label=original-54]
  \begin{tikzpicture}[remember picture,overlay,x=1bp,y=-1bp]
  \begin{scope}[shift={(current page.north west)}]
  \Fundo{54}
  \Texto{40.175}{27.803}{87.916}{17.983}{b}{2C2084}{Contagem}
  \Texto{35.433}{63.506}{310.124}{8.518}{i}{737696}{Quando todos os resultados são igualmente prováveis, probabilidade é contagem.}
  \Texto{210.648}{100.978}{66.564}{9.465}{r}{262848}{casos possíveis}
  \Texto{193.725}{163.384}{4.732}{9.465}{i}{262848}{k}
  \Texto{35.433}{186.115}{65.21}{9.465}{b}{262848}{Combinações:}
  \Texto{103.898}{186.115}{164.37}{9.465}{r}{262848}{escolher 2 alunos de uma turma de 20:}
  \Texto{277.929}{189.55}{3.684}{6.625}{r}{262848}{2}
  \Texto{49.227}{209.024}{49.711}{8.518}{r}{000000}{O coeficiente}
  \Texto{105.529}{212.067}{2.839}{5.679}{i}{000000}{k}
  \Texto{115.53}{209.024}{49.285}{8.518}{r}{000000}{conta grupos}
  \Texto{167.183}{209.024}{45.273}{8.518}{b}{000000}{sem ordem}
  \Texto{212.456}{209.024}{150.505}{8.518}{r}{000000}{; ele reaparece na distribuição binomial.}
  \Texto{37.21}{247.463}{40.499}{6.625}{b}{737696}{EMASC 2026}
  \Texto{80.653}{247.463}{2.319}{6.625}{r}{FFB014}{\textbullet{}}
  \Texto{85.913}{247.463}{117.821}{6.625}{r}{737696}{Liga Acadêmica Matemática \& Empresa}
  \Texto{206.682}{247.463}{2.319}{6.625}{r}{FFB014}{\textbullet{}}
  \Texto{211.943}{247.463}{13.25}{6.625}{r}{737696}{UFS}
  \Texto{397.549}{248.132}{7.367}{6.625}{b}{2C2084}{31}
  \Texto{406.02}{248.132}{10.315}{6.625}{r}{737696}{/ 40}
  \end{scope}
  \end{tikzpicture}
\end{frame}


% Página 55 do PDF de Gabriel
\begin{frame}[plain,noframenumbering,label=original-55]
  \begin{tikzpicture}[remember picture,overlay,x=1bp,y=-1bp]
  \begin{scope}[shift={(current page.north west)}]
  \Fundo{55}
  \Texto{40.175}{27.803}{222.48}{17.983}{b}{2C2084}{Probabilidade: definições}
  \Texto{241.22}{105.461}{5.262}{9.465}{i}{262848}{n}
  \Texto{49.227}{156.029}{357.083}{8.518}{r}{000000}{A frequência relativa estima a probabilidade de um evento. A probabilidade condicional res-}
  \Texto{49.227}{166.988}{133.792}{8.518}{r}{000000}{tringe o espaço amostral ao evento}
  \Texto{185.623}{166.988}{5.682}{8.518}{i}{000000}{B}
  \Texto{191.683}{166.988}{214.628}{8.518}{r}{000000}{; independência significa que essa restrição não altera a}
  \Texto{49.227}{177.947}{39.302}{8.518}{r}{000000}{chance de}
  \Texto{90.896}{177.947}{5.682}{8.518}{i}{000000}{A}
  \Texto{96.578}{177.947}{2.368}{8.518}{r}{000000}{.}
  \Texto{37.21}{247.463}{40.499}{6.625}{b}{737696}{EMASC 2026}
  \Texto{80.653}{247.463}{2.319}{6.625}{r}{FFB014}{\textbullet{}}
  \Texto{85.913}{247.463}{117.821}{6.625}{r}{737696}{Liga Acadêmica Matemática \& Empresa}
  \Texto{206.682}{247.463}{2.319}{6.625}{r}{FFB014}{\textbullet{}}
  \Texto{211.943}{247.463}{13.25}{6.625}{r}{737696}{UFS}
  \Texto{397.549}{248.132}{7.367}{6.625}{b}{2C2084}{32}
  \Texto{406.02}{248.132}{10.315}{6.625}{r}{737696}{/ 40}
  \end{scope}
  \end{tikzpicture}
\end{frame}


% Página 56 do PDF de Gabriel
\begin{frame}[plain,noframenumbering,label=original-56]
  \begin{tikzpicture}[remember picture,overlay,x=1bp,y=-1bp]
  \begin{scope}[shift={(current page.north west)}]
  \Fundo{56}
  \Texto{40.175}{29.691}{224.476}{17.983}{b}{2C2084}{Probabilidade condicional}
  \Texto{35.433}{63.506}{232.533}{8.518}{i}{737696}{Saber que um evento ocorreu pode alterar a chance de outro.}
  \Texto{49.227}{203.725}{357.083}{8.518}{r}{000000}{Entre os estudantes que usam mais de 6 h de redes, a proporção com pontuação de saúde}
  \Texto{49.227}{225.643}{60.112}{8.518}{b}{000000}{independentes}
  \Texto{109.338}{225.643}{2.368}{8.518}{r}{000000}{.}
  \Texto{37.21}{247.463}{40.499}{6.625}{b}{737696}{EMASC 2026}
  \Texto{80.653}{247.463}{2.319}{6.625}{r}{FFB014}{\textbullet{}}
  \Texto{85.913}{247.463}{117.821}{6.625}{r}{737696}{Liga Acadêmica Matemática \& Empresa}
  \Texto{206.682}{247.463}{2.319}{6.625}{r}{FFB014}{\textbullet{}}
  \Texto{211.943}{247.463}{13.25}{6.625}{r}{737696}{UFS}
  \Texto{397.549}{248.132}{7.367}{6.625}{b}{2C2084}{33}
  \Texto{406.02}{248.132}{10.315}{6.625}{r}{737696}{/ 40}
  \end{scope}
  \end{tikzpicture}
\end{frame}


% Páginas 57 a 58 do PDF de Gabriel
\begin{frame}[plain,noframenumbering,label=original-57]
  \begin{tikzpicture}[remember picture,overlay,x=1bp,y=-1bp]
  \begin{scope}[shift={(current page.north west)}]
  \only<1>{\Fundo{57}}
  \only<2>{\Fundo{58}}
  \Texto{40.175}{27.803}{178.512}{17.983}{b}{2C2084}{Binomial: simulação}
  \Texto{35.433}{63.506}{252.337}{8.518}{i}{737696}{Em uma turma de 20, quantos alunos esperamos com sono curto?}
  \only<1>{\Texto{49.227}{209.198}{357.083}{8.518}{r}{000000}{Turmas de 20 estudantes sorteadas dos dados; conta-se quantos dormem menos de 6 h,}}
  \only<1>{\Texto{49.227}{220.157}{190.198}{8.518}{r}{000000}{evento que ocorre em cerca de 33\% dos registros.}}
  \Texto{37.21}{247.463}{40.499}{6.625}{b}{737696}{EMASC 2026}
  \Texto{80.653}{247.463}{2.319}{6.625}{r}{FFB014}{\textbullet{}}
  \Texto{85.913}{247.463}{117.821}{6.625}{r}{737696}{Liga Acadêmica Matemática \& Empresa}
  \Texto{206.682}{247.463}{2.319}{6.625}{r}{FFB014}{\textbullet{}}
  \Texto{211.943}{247.463}{13.25}{6.625}{r}{737696}{UFS}
  \Texto{397.549}{248.132}{7.367}{6.625}{b}{2C2084}{34}
  \Texto{406.02}{248.132}{10.315}{6.625}{r}{737696}{/ 40}
  \only<2>{\Texto{49.227}{209.394}{357.083}{8.518}{r}{000000}{Repetindo o sorteio 10.000 vezes, as frequências observadas aproximam-se das probabilida-}}
  \only<2>{\Texto{49.227}{220.353}{25.571}{8.518}{r}{000000}{des da}}
  \only<2>{\Texto{77.165}{220.353}{85.504}{8.518}{b}{000000}{distribuição binomial}}
  \only<2>{\Texto{162.668}{220.353}{2.368}{8.518}{r}{000000}{.}}
  \end{scope}
  \end{tikzpicture}
\end{frame}


% Página 59 do PDF de Gabriel
\begin{frame}[plain,noframenumbering,label=original-59]
  \begin{tikzpicture}[remember picture,overlay,x=1bp,y=-1bp]
  \begin{scope}[shift={(current page.north west)}]
  \Fundo{59}
  \Texto{40.175}{29.916}{156.502}{17.983}{b}{2C2084}{Binomial: fórmula}
  \Texto{109.687}{80.832}{4.732}{9.465}{i}{262848}{k}
  \Texto{236.631}{87.561}{55.553}{6.974}{t}{000000}{n,\ p\ =\ 20,\ 0.33}
  \Texto{236.631}{105.494}{88.881}{6.974}{t}{000000}{stats.binom.pmf(0,\ n,\ p)}
  \Texto{336.637}{105.494}{29.625}{6.974}{c}{737696}{\#\ P(X=0)}
  \Texto{236.631}{114.46}{85.178}{6.974}{t}{000000}{stats.binom.sf(9,\ n,\ p)}
  \Texto{336.637}{114.46}{37.031}{6.974}{c}{737696}{\#\ P(X>=10)}
  \Texto{236.631}{123.427}{51.843}{6.974}{t}{000000}{n*p,\ n*p*(1-p)}
  \Texto{336.637}{123.427}{66.663}{6.974}{c}{737696}{\#\ media,\ variancia}
  \Texto{49.227}{169.395}{32.952}{8.966}{t}{2C2084}{sf(k-1)}
  \Texto{82.17}{169.395}{103.306}{8.518}{r}{000000}{. O modelo exige tentativas}
  \Texto{187.843}{169.395}{60.112}{8.518}{b}{000000}{independentes}
  \Texto{250.322}{169.395}{58.238}{8.518}{r}{000000}{e probabilidade}
  \Texto{310.927}{169.395}{40.231}{8.518}{b}{000000}{constante}
  \Texto{351.157}{169.395}{2.368}{8.518}{r}{000000}{.}
  \Texto{37.21}{247.463}{40.499}{6.625}{b}{737696}{EMASC 2026}
  \Texto{80.653}{247.463}{2.319}{6.625}{r}{FFB014}{\textbullet{}}
  \Texto{85.913}{247.463}{117.821}{6.625}{r}{737696}{Liga Acadêmica Matemática \& Empresa}
  \Texto{206.682}{247.463}{2.319}{6.625}{r}{FFB014}{\textbullet{}}
  \Texto{211.943}{247.463}{13.25}{6.625}{r}{737696}{UFS}
  \Texto{397.549}{248.132}{7.367}{6.625}{b}{2C2084}{35}
  \Texto{406.02}{248.132}{10.315}{6.625}{r}{737696}{/ 40}
  \end{scope}
  \end{tikzpicture}
\end{frame}


% Página 60 do PDF de Gabriel
\begin{frame}[plain,noframenumbering,label=original-60]
  \begin{tikzpicture}[remember picture,overlay,x=1bp,y=-1bp]
  \begin{scope}[shift={(current page.north west)}]
  \Fundo{60}
  \Texto{40.175}{27.803}{171.805}{17.983}{b}{2C2084}{Poisson: simulação}
  \Texto{35.433}{63.506}{275.421}{8.518}{i}{737696}{Para eventos raros em grupos grandes, contamos ocorrências por grupo.}
  \Texto{49.227}{203.827}{298.454}{8.518}{r}{000000}{Grupos de 100 estudantes; conta-se quantos dormem menos de 4 h, em média}
  \Texto{349.87}{203.827}{54.072}{8.518}{b}{000000}{2,3 por grupo}
  \Texto{403.941}{203.827}{2.368}{8.518}{r}{000000}{.}
  \Texto{49.227}{214.786}{260.429}{8.518}{r}{000000}{A curva corresponde à distribuição de Poisson com essa taxa média.}
  \Texto{37.21}{247.463}{40.499}{6.625}{b}{737696}{EMASC 2026}
  \Texto{80.653}{247.463}{2.319}{6.625}{r}{FFB014}{\textbullet{}}
  \Texto{85.913}{247.463}{117.821}{6.625}{r}{737696}{Liga Acadêmica Matemática \& Empresa}
  \Texto{206.682}{247.463}{2.319}{6.625}{r}{FFB014}{\textbullet{}}
  \Texto{211.943}{247.463}{13.25}{6.625}{r}{737696}{UFS}
  \Texto{397.549}{248.132}{7.367}{6.625}{b}{2C2084}{36}
  \Texto{406.02}{248.132}{10.315}{6.625}{r}{737696}{/ 40}
  \end{scope}
  \end{tikzpicture}
\end{frame}


% Página 61 do PDF de Gabriel
\begin{frame}[plain,noframenumbering,label=original-61]
  \begin{tikzpicture}[remember picture,overlay,x=1bp,y=-1bp]
  \begin{scope}[shift={(current page.north west)}]
  \Fundo{61}
  \Texto{40.175}{29.916}{149.794}{17.983}{b}{2C2084}{Poisson: fórmula}
  \Texto{236.631}{87.561}{33.328}{6.974}{t}{000000}{lam\ =\ 2.3}
  \Texto{236.631}{105.494}{92.584}{6.974}{t}{000000}{stats.poisson.pmf(0,\ lam)}
  \Texto{340.341}{105.494}{59.256}{6.974}{c}{737696}{\#\ P(X=0)\ \textasciitilde{}\ 0.100}
  \Texto{236.631}{114.46}{88.881}{6.974}{t}{000000}{stats.poisson.sf(4,\ lam)}
  \Texto{340.341}{114.46}{62.96}{6.974}{c}{737696}{\#\ P(X>=5)\ \textasciitilde{}\ 0.084}
  \Texto{49.227}{149.01}{217.541}{8.518}{r}{000000}{A distribuição de Poisson é o limite da binomial quando}
  \Texto{270.164}{149.01}{4.736}{8.518}{i}{000000}{n}
  \Texto{278.434}{149.01}{27.326}{8.518}{r}{000000}{cresce,}
  \Texto{309.409}{149.01}{4.736}{8.518}{i}{000000}{p}
  \Texto{317.777}{149.01}{75.427}{8.518}{r}{000000}{diminui e o produto}
  \Texto{396.599}{149.01}{9.472}{8.518}{i}{000000}{np}
  \Texto{49.227}{159.969}{309.289}{8.518}{r}{000000}{permanece constante. Nas contagens por grupo, média e variância são próximas.}
  \Texto{37.21}{247.463}{40.499}{6.625}{b}{737696}{EMASC 2026}
  \Texto{80.653}{247.463}{2.319}{6.625}{r}{FFB014}{\textbullet{}}
  \Texto{85.913}{247.463}{117.821}{6.625}{r}{737696}{Liga Acadêmica Matemática \& Empresa}
  \Texto{206.682}{247.463}{2.319}{6.625}{r}{FFB014}{\textbullet{}}
  \Texto{211.943}{247.463}{13.25}{6.625}{r}{737696}{UFS}
  \Texto{397.549}{248.132}{7.367}{6.625}{b}{2C2084}{37}
  \Texto{406.02}{248.132}{10.315}{6.625}{r}{737696}{/ 40}
  \end{scope}
  \end{tikzpicture}
\end{frame}


% Página 62 do PDF de Gabriel
\begin{frame}[plain,noframenumbering,label=original-62]
  \begin{tikzpicture}[remember picture,overlay,x=1bp,y=-1bp]
  \begin{scope}[shift={(current page.north west)}]
  \Fundo{62}
  \Texto{40.175}{29.691}{179.142}{17.983}{b}{2C2084}{Binomial ou Poisson}
  \Texto{35.433}{63.506}{245.83}{8.518}{i}{737696}{Dois modelos de contagem: a escolha depende do que se conta.}
  \Texto{38.534}{100.643}{39.969}{9.465}{b}{027499}{Binomial}
  \Texto{38.534}{115.586}{25.251}{9.465}{r}{000000}{Conta}
  \Texto{66.416}{115.586}{59.447}{9.465}{b}{000000}{sucessos em}
  \Texto{128.493}{115.586}{5.262}{9.465}{i}{000000}{n}
  \Texto{136.542}{115.586}{43.991}{9.465}{b}{000000}{tentativas}
  \Texto{180.532}{115.586}{2.631}{9.465}{r}{000000}{.}
  \Texto{38.534}{127.542}{65.116}{9.465}{r}{000000}{Limite superior:}
  \Texto{106.904}{127.542}{5.262}{9.465}{i}{000000}{n}
  \Texto{112.322}{127.542}{2.631}{9.465}{r}{000000}{.}
  \Texto{38.534}{171.377}{158.767}{8.518}{r}{737696}{Ex.: estudantes com sono curto em turma}
  \Texto{38.534}{183.332}{23.68}{8.518}{r}{737696}{de 20.}
  \Texto{250.943}{100.642}{36.438}{9.465}{b}{CC8D10}{Poisson}
  \Texto{250.943}{115.586}{25.251}{9.465}{r}{000000}{Conta}
  \Texto{278.825}{115.586}{128.774}{9.465}{b}{000000}{ocorrências em um intervalo}
  \Texto{407.597}{115.586}{2.631}{9.465}{r}{000000}{.}
  \Texto{250.943}{127.542}{83.288}{9.465}{r}{000000}{Sem limite superior.}
  \Texto{250.943}{159.422}{154.253}{8.518}{r}{737696}{Ex.: chamadas por hora em uma central.}
  \Texto{49.227}{213.169}{90.904}{8.518}{r}{000000}{Ambas supõem eventos}
  \Texto{142.498}{213.169}{60.112}{8.518}{b}{000000}{independentes}
  \Texto{204.977}{213.169}{64.745}{8.518}{r}{000000}{e taxa constante.}
  \Texto{37.21}{247.463}{40.499}{6.625}{b}{737696}{EMASC 2026}
  \Texto{80.653}{247.463}{2.319}{6.625}{r}{FFB014}{\textbullet{}}
  \Texto{85.913}{247.463}{117.821}{6.625}{r}{737696}{Liga Acadêmica Matemática \& Empresa}
  \Texto{206.682}{247.463}{2.319}{6.625}{r}{FFB014}{\textbullet{}}
  \Texto{211.943}{247.463}{13.25}{6.625}{r}{737696}{UFS}
  \Texto{397.549}{248.132}{7.367}{6.625}{b}{2C2084}{38}
  \Texto{406.02}{248.132}{10.315}{6.625}{r}{737696}{/ 40}
  \end{scope}
  \end{tikzpicture}
\end{frame}


\section{Síntese da análise}

% Página 63 do PDF de Gabriel
\begin{frame}[plain,noframenumbering,label=original-63]
  \begin{tikzpicture}[remember picture,overlay,x=1bp,y=-1bp]
  \begin{scope}[shift={(current page.north west)}]
  \Fundo{63}
  \Texto{44.456}{41.894}{111.824}{8.518}{b}{34E0C8}{SÍNTESE E CONTINUIDADE}
  \Texto{43.039}{116.125}{63.147}{56.787}{b}{FFB014}{08}
  \Texto{43.039}{155.148}{215.44}{24.608}{b}{FFFFFF}{Síntese da análise}
  \Texto{43.039}{234.662}{46.285}{7.572}{b}{34E0C8}{EMASC 2026}
  \Texto{91.429}{234.662}{2.65}{7.572}{b}{FFB014}{\textbullet{}}
  \Texto{96.176}{234.662}{21.45}{7.572}{b}{34E0C8}{LAME}
  \Link{42.042}{135.586}{217.431}{20.866}{original-63}
  \end{scope}
  \end{tikzpicture}
\end{frame}


% Página 64 do PDF de Gabriel
\begin{frame}[plain,noframenumbering,label=original-64]
  \begin{tikzpicture}[remember picture,overlay,x=1bp,y=-1bp]
  \begin{scope}[shift={(current page.north west)}]
  \Fundo{64}
  \Texto{40.175}{27.942}{190.525}{17.983}{b}{2C2084}{Principais conclusões}
  \Texto{35.433}{63.506}{175.76}{8.518}{i}{737696}{Quatro práticas resumem o percurso do curso.}
  \Texto{28.574}{100.642}{42.145}{9.465}{b}{027499}{1. Validar}
  \Texto{28.574}{115.586}{67.87}{9.465}{r}{000000}{A qualidade dos}
  \Texto{28.574}{127.542}{67.349}{9.465}{r}{000000}{dados antecede}
  \Texto{28.574}{139.497}{39.599}{9.465}{r}{000000}{o cálculo.}
  \Texto{137.268}{100.642}{56.768}{9.465}{b}{007078}{2. Dispersão}
  \Texto{137.268}{115.586}{67.331}{9.465}{r}{000000}{A média isolada}
  \Texto{137.268}{127.542}{57.866}{9.465}{r}{000000}{é insuficiente:}
  \Texto{137.268}{139.497}{44.19}{9.465}{r}{000000}{mediana e}
  \Texto{137.268}{151.452}{41.029}{9.465}{r}{000000}{dispersão}
  \Texto{137.268}{163.407}{54.175}{9.465}{r}{000000}{completam o}
  \Texto{137.268}{175.362}{33.807}{9.465}{r}{000000}{resumo.}
  \Texto{245.963}{100.879}{44.284}{9.465}{b}{231A6A}{3. Gráfico}
  \Texto{245.963}{115.823}{38.303}{9.465}{r}{000000}{O gráfico}
  \Texto{245.963}{127.778}{46.83}{9.465}{r}{000000}{antecede o}
  \Texto{245.963}{139.733}{47.729}{9.465}{r}{000000}{coeficiente.}
  \Texto{354.657}{100.642}{39.552}{9.465}{b}{CC8D10}{4. Causa}
  \Texto{354.657}{115.586}{66.81}{9.465}{r}{000000}{Associação não}
  \Texto{354.657}{127.541}{29.444}{9.465}{r}{000000}{implica}
  \Texto{354.657}{139.497}{52.992}{9.465}{r}{000000}{causalidade.}
  \Texto{37.21}{247.463}{40.499}{6.625}{b}{737696}{EMASC 2026}
  \Texto{80.653}{247.463}{2.319}{6.625}{r}{FFB014}{\textbullet{}}
  \Texto{85.914}{247.463}{117.821}{6.625}{r}{737696}{Liga Acadêmica Matemática \& Empresa}
  \Texto{206.683}{247.463}{2.319}{6.625}{r}{FFB014}{\textbullet{}}
  \Texto{211.943}{247.463}{13.25}{6.625}{r}{737696}{UFS}
  \Texto{397.549}{248.132}{7.367}{6.625}{b}{2C2084}{39}
  \Texto{406.02}{248.132}{10.315}{6.625}{r}{737696}{/ 40}
  \end{scope}
  \end{tikzpicture}
\end{frame}


% Página 65 do PDF de Gabriel
\begin{frame}[plain,noframenumbering,label=original-65]
  \begin{tikzpicture}[remember picture,overlay,x=1bp,y=-1bp]
  \begin{scope}[shift={(current page.north west)}]
  \Fundo{65}
  \Texto{40.175}{28.167}{149.902}{17.983}{b}{2C2084}{Próximos passos}
  \Texto{35.433}{63.506}{215.361}{8.518}{i}{737696}{Descrever os dados abre caminho para estimar e decidir.}
  \Texto{38.534}{100.879}{75.555}{9.465}{b}{231A6A}{Análise preditiva}
  \Texto{38.534}{115.823}{137.481}{9.465}{r}{000000}{Estimar valores não observados.}
  \Texto{38.534}{131.763}{162.975}{8.518}{r}{737696}{Ex.: prever a pontuação de saúde mental a}
  \Texto{38.534}{143.718}{122.881}{8.518}{r}{737696}{partir das horas de uso de redes}
  \Texto{38.534}{155.674}{45.827}{8.518}{r}{737696}{(regressão).}
  \Texto{250.943}{100.879}{83.979}{9.465}{b}{CC8D10}{Análise prescritiva}
  \Texto{250.943}{115.823}{132.408}{9.465}{r}{000000}{Recomendar ações a partir das}
  \Texto{250.943}{127.778}{49.584}{9.465}{r}{000000}{estimativas.}
  \Texto{250.943}{143.718}{147.319}{8.518}{r}{737696}{Ex.: definir limites de uso associados a}
  \Texto{250.943}{155.674}{113.119}{8.518}{r}{737696}{maior bem-estar (otimização).}
  \Texto{49.227}{187.282}{294.28}{8.518}{r}{000000}{Ambas exigem dados confiáveis e, no caso de recomendações, evidência de}
  \Texto{346.268}{187.282}{48.766}{8.518}{b}{000000}{causalidade}
  \Texto{397.791}{187.282}{8.518}{8.518}{r}{000000}{\textemdash{}}
  \Texto{49.227}{198.241}{91.517}{8.518}{r}{000000}{não apenas associação.}
  \Texto{35.433}{217.652}{168.079}{6.625}{r}{737696}{Para continuar: documentação do pandas e do seaborn \textbullet{}}
  \Texto{205.352}{217.652}{109.082}{6.625}{i}{737696}{Practical Statistics for Data Scientists}
  \Texto{316.467}{217.652}{47.323}{6.625}{r}{737696}{(Bruce \& Bruce)}
  \Texto{37.21}{247.463}{40.499}{6.625}{b}{737696}{EMASC 2026}
  \Texto{80.653}{247.463}{2.319}{6.625}{r}{FFB014}{\textbullet{}}
  \Texto{85.913}{247.463}{117.821}{6.625}{r}{737696}{Liga Acadêmica Matemática \& Empresa}
  \Texto{206.682}{247.463}{2.319}{6.625}{r}{FFB014}{\textbullet{}}
  \Texto{211.943}{247.463}{13.25}{6.625}{r}{737696}{UFS}
  \Texto{397.549}{248.132}{7.367}{6.625}{b}{2C2084}{40}
  \Texto{406.02}{248.132}{10.315}{6.625}{r}{737696}{/ 40}
  \end{scope}
  \end{tikzpicture}
\end{frame}


% Página 66 do PDF de Gabriel
\begin{frame}[plain,noframenumbering,label=original-66]
  \begin{tikzpicture}[remember picture,overlay,x=1bp,y=-1bp]
  \begin{scope}[shift={(current page.north west)}]
  \Fundo{66}
  \Texto{38.787}{87.092}{198.974}{41.644}{b}{FFFFFF}{Obrigado!}
  \Texto{38.787}{128.21}{76.68}{17.036}{b}{34E0C8}{Dúvidas?}
  \Texto{38.787}{154.186}{64.728}{8.518}{r}{D1CEE4}{Gabriel C. Souza}
  \Texto{38.787}{166.141}{177.992}{8.518}{r}{D1CEE4}{Liga Acadêmica Matemática \& Empresa \textendash{} UFS}
  \Texto{38.787}{178.096}{96.816}{8.518}{r}{D1CEE4}{email-da-liga@dominio.br}
  \end{scope}
  \end{tikzpicture}
\end{frame}



\end{document}
